\pdfoutput=1

\documentclass[11pt]{article}
\usepackage{algorithm}
\usepackage{algpseudocode}
\usepackage{colortbl}  % 배경색
\usepackage{xcolor}    % 색상
\usepackage{pifont}    % 체크/X 기호
\usepackage{makecell}  % 셀 내 줄바꿈 기능 (\makecell)
\definecolor{GrayHeader}{gray}{0.9}
\definecolor{CheckGreen}{rgb}{0.1, 0.6, 0.1}
\definecolor{CrossRed}{rgb}{0.8, 0.0, 0.0}
\newcommand{\cmark}{\textcolor{CheckGreen}{\ding{51}}}
\newcommand{\xmark}{\textcolor{CrossRed}{\ding{55}}}
\usepackage[final]{acl}
\usepackage{kotex}
\usepackage{times}
\usepackage{latexsym}
\usepackage[T1]{fontenc}
% \usepackage[T5]{fontenc}
\usepackage[utf8]{inputenc}
\usepackage{microtype}
\usepackage{inconsolata}
\usepackage{graphicx}
\usepackage{enumitem}
\usepackage{booktabs}  % 모던 테이블 스타일
\usepackage{multirow}
\usepackage{adjustbox}  % 테이블 크기 조정
\usepackage{listings}   % 코드 블록
\definecolor{LightBlue}{RGB}{230, 242, 255}  % 하이라이트 배경색
\definecolor{CodeBg}{RGB}{245, 245, 245}     % 코드 배경색

% S.O.L.I.D. 컴포넌트별 색상
\definecolor{ShiftColor}{RGB}{66, 133, 244}    % S - Blue
\definecolor{OracleColor}{RGB}{52, 168, 83}    % O - Green
\definecolor{LiveColor}{RGB}{251, 188, 5}      % L - Yellow/Gold
\definecolor{InterColor}{RGB}{234, 67, 53}     % I - Red
\definecolor{DefensiveColor}{RGB}{156, 39, 176}% D - Purple
\definecolor{SafetyColor}{RGB}{96, 125, 139}   % Safety - Blue Gray

% listings 설정
\lstset{
    basicstyle=\ttfamily\small,
    backgroundcolor=\color{CodeBg},
    frame=single,
    framesep=3pt,
    breaklines=true,
    breakatwhitespace=true,
    columns=fullflexible,
    keepspaces=true,
    showstringspaces=false,
    xleftmargin=5pt,
    xrightmargin=5pt,
}

% tcolorbox 설정
\usepackage[most]{tcolorbox}

% 프롬프트 박스 스타일 (컴포넌트별)
\newtcolorbox{promptbox}[2][]{%
    enhanced,
    colback=#2!5!white,
    colframe=#2!75!black,
    fonttitle=\bfseries\sffamily,
    attach boxed title to top left={yshift=-2mm, xshift=3mm},
    boxed title style={colback=#2!75!black, colframe=#2!75!black},
    top=4mm,
    boxrule=0.5pt,
    #1
}

% 코드 박스 스타일
\newtcolorbox{codebox}[1][]{%
    enhanced,
    colback=CodeBg,
    colframe=gray!50!black,
    fonttitle=\bfseries\sffamily,
    attach boxed title to top left={yshift=-2mm, xshift=3mm},
    boxed title style={colback=gray!60, colframe=gray!50!black},
    top=4mm,
    boxrule=0.4pt,
    left=2mm, right=2mm,
    #1
}

% Safety 메커니즘 박스
\newtcolorbox{safetybox}[1][]{%
    enhanced,
    colback=SafetyColor!8!white,
    colframe=SafetyColor!80!black,
    fonttitle=\bfseries\sffamily,
    attach boxed title to top left={yshift=-2mm, xshift=3mm},
    boxed title style={colback=SafetyColor!80!black, colframe=SafetyColor!80!black},
    top=4mm,
    boxrule=0.5pt,
    #1
}

% 전체 예시 박스 (breakable, 여러 페이지에 걸침)
\newtcolorbox{examplebox}[1][]{%
    enhanced,
    breakable,
    colback=gray!3!white,
    colframe=black!70,
    fonttitle=\bfseries\sffamily\large,
    attach boxed title to top center={yshift=-3mm},
    boxed title style={colback=black!70, colframe=black!70},
    top=6mm,
    boxrule=0.6pt,
    left=3mm, right=3mm,
    before upper={\parindent0pt\small},
    #1
}

% 코드 리스팅 박스 (listings 통합)
\newtcblisting{codelistbox}[2][]{%
    enhanced,
    colback=#2!5!white,
    colframe=#2!80!black,
    fonttitle=\bfseries\sffamily,
    attach boxed title to top left={yshift=-2mm, xshift=3mm},
    boxed title style={colback=#2!80!black, colframe=#2!80!black},
    top=2mm, bottom=2mm,
    boxrule=0.4pt,
    left=2mm, right=2mm,
    listing only,
    listing options={
        language=Python,
        basicstyle=\ttfamily\footnotesize,
        breaklines=true,
        columns=flexible,
        keepspaces=true,
        showstringspaces=false,
        commentstyle=\color{gray},
        keywordstyle=\color{blue!70!black},
        stringstyle=\color{orange!80!black},
        morekeywords={self, True, False, None},
    },
    #1
}

\title{SolidCoder: Bridging the Mental-Reality Gap in LLM Code Generation through Concrete Execution}

\author{
Woojin Lee \quad
Jin-Xia Huang \\
\\Electronics and Telecommunications Research Institute, Republic of Korea  \\
\texttt{\{writerwoody, hgh\}@etri.re.kr} \\
} 

\begin{document}
\maketitle
\begin{abstract}
% 대형 언어 모델(LLMs)은 자동화된 코드 생성에서 놀라운 진전을 이끌어냈습니다. CodeSIM과 같은 최첨단 멀티-에이전트 프레임워크는 \textit{정신 시뮬레이션}---모델이 코드 실행을 내부적으로 추적하는 곳---을 활용하여 계획을 검증하고 오류를 디버깅합니다. 그러나 이 접근법은 \textbf{정신-현실 격차}에 의해 근본적으로 제한됩니다: LLM들은 실행 추적을 자주 환각하여, 잘못된 코드를 올바른 것으로 믿습니다. 우리는 \textbf{SolidCoder}를 소개합니다. 이 프레임워크는 검증을 \textbf{구체적 현실}---상상된 추적 대신 실제 코드 실행---에 고정함으로써 이 격차를 연결합니다. 우리의 \textbf{S.O.L.I.D.} 아키텍처는 다섯 가지 시너지적 구성 요소를 통합합니다: \textbf{S}hift-left 엣지-케이스 계획, \textbf{O}racle 기반 속성 단언, \textbf{L}ive 샌드박스드 실행, \textbf{I}ntermediate 코드 시뮬레이션, 그리고 \textbf{D}efensive 테스트 축적. GPT-4o를 사용한 실험들은 SolidCoder가 이전 최첨단인 CodeSIM을 지속적으로 능가한다는 것을 보여주며, 새로운 pass@1 결과를 달성합니다: HumanEval에서 \textbf{95.7\%} (CodeSIM 대비 +0.6\%p), CodeContests에서 \textbf{77.0\%} (+4.3\%p), 그리고 APPS에서 \textbf{26.7\%} (+3.4\%p). 이러한 개선은 RL 기반 post-training을 거친 최신 모델들에서도 일관되게 나타납니다. 이는 에이전트를 런타임 현실에 기반짓는 것이 견고한 코드 합성에 필수적임을 검증합니다. 우리는 미래 연구를 촉진하기 위해 우리의 코드와 실험 프레임워크를 공개합니다.\footnote{\url{https://github.com/Anonymous-ARR-Submissions/SolidCoder}}
Large Language Models (LLMs) have driven remarkable progress in automated code generation. State-of-the-art multi-agent frameworks such as CodeSIM leverage \textit{Mental Simulation}---where the model internally traces code execution---to verify plans and debug errors. However, this approach is fundamentally limited by the \textbf{Mental-Reality Gap}: LLMs frequently hallucinate execution traces, believing incorrect code to be correct. We introduce \textbf{SolidCoder}, a framework that bridges this gap by anchoring verification in \textbf{Concrete Reality}---actual code execution rather than imagined traces. Our \textbf{S.O.L.I.D.} architecture integrates five synergistic components: \textbf{S}hift-left edge-case planning, \textbf{O}racle-based property assertions, \textbf{L}ive sandboxed execution, \textbf{I}ntermediate code simulation, and \textbf{D}efensive test accumulation. Experiments with GPT-4o demonstrate that SolidCoder consistently outperforms the previous state-of-the-art, CodeSIM, achieving new pass@1 results: \textbf{95.7\%} on HumanEval (+0.6\%p over CodeSIM), \textbf{77.0\%} on CodeContests (+4.3\%p), and \textbf{26.7\%} on APPS (+3.4\%p). These improvements generalize consistently across recent RL post-trained models. This validates that grounding agents in runtime reality is essential for robust code synthesis. We release our code and experimental framework to facilitate future research.\footnote{\url{https://github.com/Anonymous-ARR-Submissions/SolidCoder}}
\end{abstract}

\section{Introduction}

\begin{figure*}[t]
    \centering
    \includegraphics[width=\textwidth]{fig1.png}
    % \caption{\textbf{비교 아키텍처 개요: CodeSIM 대 SolidCoder} 상단 파이프라인은 CodeSIM 기준선을 보여줍니다. CodeSIM 기준선은 계획 및 디버깅 중 검증을 위해 정신적 시뮬레이션("상상")에 의존하며, 이로 인해 LLM 환각에 취약한 "정신-현실 간극"이 발생합니다. 하단 파이프라인은 SolidCoder를 보여주는데, SolidCoder는 검증을 구체적인 실행("실행")에 기반하여 이러한 간극을 메웁니다. SolidCoder는 \textcolor{blue}{blue}로 강조된 S.O.L.I.D. 개선 사항을 통합합니다. (1) \textbf{Planning Agent}는 \textbf{S}hift-left를 통합하여 계획 생성 *전에* 경계 사례를 식별합니다. (2) \textbf{Coding Agent}는 \textbf{I}ntermediate Simulation을 적용하여 코드 추적을 검증합니다. (3) 중앙 라이브 검증 루프(Live Verification Loop)는 오라클 기반 어설션을 활용하는 샌드박스 라이브 실행(Sandbox Live Execution)으로 멘탈 체크를 대체합니다. 실패한 테스트는 방어적 누적(Defensive Accumulation)을 통해 저장되며, 수정된 코드는 회귀를 방지하기 위해 모든 누적 테스트([D])를 통과해야 솔루션으로 채택됩니다.}
    \caption{\textbf{Comparative Architecture Overview: CodeSIM vs. SolidCoder.} The top pipeline illustrates the CodeSIM baseline, which relies on mental simulation ("imagines") for verification during planning and debugging, resulting in a "Mental-Reality Gap" prone to LLM hallucination. The bottom pipeline demonstrates SolidCoder, which bridges this gap by grounding verification in concrete execution ("executes"). SolidCoder incorporates S.O.L.I.D. enhancements highlighted in \textcolor{blue}{blue}: (1) The \textbf{Planning Agent} integrates \textbf{S}hift-left to identify edge cases before plan generation; (2) The \textbf{Coding Agent} applies \textbf{I}ntermediate Simulation to verify code traces; (3) The central \textbf{Live Verification Loop} replaces mental checks with sandbox \textbf{L}ive Execution utilizing \textbf{O}racle-based assertions. Failing tests are stored via \textbf{D}efensive Accumulation, and fixed code must pass *all* accumulated tests ([D]) to prevent regression before accepted as a solution.}
    \label{fig:1}
\end{figure*}

% ============================================================
% [Figure 1 Drawing Instructions - SolidCoder Architecture]
% ============================================================
%
% 전체 레이아웃: 좌→우 흐름의 3단계 파이프라인 (CodeSIM과 유사하되 S.O.L.I.D. 강조)
% 크기: 전체 텍스트 너비 (\textwidth), 세로 약 6-7cm
%
% ========== 왼쪽: INPUT ==========
% [Problem Box]
%   - 라벨: "Problem"
%   - 내용: 문제 설명 텍스트 (예: "Given n, return the nth Fibonacci number...")
%   - 아래에 "Sample I/O" 박스 (입출력 예시)
%
% ========== 1단계: Planning Agent ==========
% [Planning Agent 큰 박스 - 연한 회색 배경]
%   라벨: "Planning Agent"
%
%   내부 흐름 (위→아래 또는 좌→우):
%
%   (1) [S] Shift-left Edge Case Identification (파란색 강조 박스)
%       - 아이콘: ⚠️ 또는 경고 삼각형
%       - 내용: "Identify edge cases:\n- Empty input\n- Boundary values\n- Large n overflow"
%       - CodeSIM에는 없는 NEW 컴포넌트임을 강조
%
%   (2) Generate Exemplar & Plan (기존 CodeSIM과 동일)
%       - 아이콘: �� 또는 문서
%       - 내용: "Retrieve similar problems\nGenerate algorithmic plan"
%
%   (3) Simulate & Verify Plan (기존 Mental Simulation)
%       - 아이콘: �� 또는 두뇌
%       - 내용: "Mental simulation\nto verify plan"
%       - 화살표: 실패 시 → Revise Plan (루프)
%
%   하단에 루프 화살표: "p times" (최대 5회 반복)
%
% ========== 2단계: Coding Agent ==========
% [Coding Agent 큰 박스 - 연한 회색 배경]
%   라벨: "Coding Agent"
%
%   내부 흐름:
%
%   (1) Generate Code (기존과 동일)
%       - 아이콘: �� 또는 코드 블록
%       - 내용: "Plan → Code\nTranslation"
%
%   (2) [I] Intermediate Simulation (파란색 강조 박스)
%       - 아이콘: �� 또는 돋보기
%       - 내용: "Trace code on\nsample input"
%       - 화살표: 불일치 시 → 다시 Generate Code
%       - 목적: Plan과 Code 간 변환 오류 조기 포착
%
% ========== 3단계: Live Verification Loop (NEW - 전체가 파란색 테두리) ==========
% [Live Verification 큰 박스 - 파란색 테두리 강조]
%   라벨: "Live Verification Loop" (파란색)
%
%   내부 흐름 (순환 구조):
%
%   (1) [O] Oracle-based Test Generation (파란색 박스)
%       - 아이콘: �� 또는 과녁
%       - 내용: "Generate property-based\nassertions without\nground-truth"
%
%   (2) [L] Live Execution (파란색 박스)
%       - 아이콘: ▶️ 또는 실행 버튼
%       - 내용: "Execute in sandbox\n(5s timeout)"
%       - 화살표 두 갈래:
%         - PASS → "Check Sample I/O" → 성공 시 OUTPUT으로
%         - FAIL → Debugging Agent로
%
%   (3) [D] Defensive Test Accumulation (파란색 박스)
%       - 아이콘: ��️ 또는 방패
%       - 내용: "Accumulate failed\ntest cases"
%       - 테스트 목록 아이콘 (체크박스들)
%
% ========== 4단계: Debugging Agent ==========
% [Debugging Agent 박스 - 연한 회색 배경]
%   라벨: "Debugging Agent"
%
%   내부 흐름:
%
%   (1) Analyze Failure
%       - 실패 원인 분석
%       - Error message + 실패한 테스트 케이스 표시
%
%   (2) Fix Code
%       - 코드 수정
%       - 화살표: 수정된 코드 → Live Verification Loop로 복귀
%       - "[D] Re-run ALL accumulated tests" 강조
%
%   하단에 루프 화살표: "d times" (최대 5회 반복)
%
% ========== 오른쪽: OUTPUT ==========
% [Solution Box]
%   - 라벨: "Solution"
%   - 내용: 최종 코드 (Python 코드 스니펫)
%   - 체크마크 ✓ 표시
%
% ========== 화살표 및 흐름 ==========
% - 메인 흐름: 굵은 검은색 화살표 (Problem → Planning → Coding → Live Verification → Solution)
% - 루프 흐름: 점선 또는 회색 화살표 (실패 시 되돌아가는 경로)
% - S.O.L.I.D. 컴포넌트: 파란색으로 강조, [S], [O], [L], [I], [D] 라벨 명시
%
% ========== 색상 팔레트 ==========
% - 기존 컴포넌트 (CodeSIM 공유): 연한 회색 배경 (#F5F5F5)
% - S.O.L.I.D. 신규 컴포넌트: 파란색 강조 (#4A90D9 또는 유사)
% - 화살표: 검은색 (메인), 회색 (루프)
% - 텍스트: 검은색, 컴포넌트 라벨은 볼드
%
% ========== 범례 (Legend) - 하단 또는 우측 ==========
% [기존 컴포넌트] - 회색 박스
% [S.O.L.I.D. 컴포넌트] - 파란색 박스
%
% ========== CodeSIM Figure 1과의 차이점 ==========
% 1. Shift-left Planning [S]: Planning Agent 첫 단계에 추가
% 2. Intermediate Simulation [I]: Coding Agent에 추가
% 3. Live Verification Loop: 전체 새로운 단계
%    - Oracle-based Assertions [O]
%    - Live Execution [L]
%    - Defensive Accumulation [D]
% 4. Debugging Agent: accumulated tests 재실행 강조
%
% ========== 대안 레이아웃 (간소화 버전) ==========
% 만약 공간이 부족하면, 상/하 2단 구조로:
% - 상단: CodeSIM 파이프라인 (회색, "Mental Simulation" 라벨)
% - 하단: SolidCoder 파이프라인 (파란색 강조, "Concrete Execution" 라벨)
% - 비교를 통해 Mental-Reality Gap을 시각적으로 표현
%
% ============================================================

% 최근 대규모 언어 모델(LLM)은 놀라운 프로그래밍 성능을 보여주었습니다. GPT, Claude, Gemini, Kimi, DeepSeek, Qwen과 같은 파운데이션 모델 시리즈들이 자연어로부터 코드를 생성하는 데 있어 인상적인 능력을 보여주고 있습니다. 그러나 정확한 프로그램을 합성하는 것은 본질적으로 어려운 과제이며, 자연어 처리, 알고리즘 설계, 데이터 구조 선택 및 문제 해결 전략에 대한 깊은 이해가 필요합니다. 이러한 과제는 엄격한 제약 조건을 준수하고 포괄적인 테스트 스위트를 통과하는 것이 매우 중요한 경쟁 프로그래밍 및 고급 소프트웨어 엔지니어링 작업의 경우 더욱 복잡해집니다 \citep{chen2021evaluating, hendrycks2021measuring, li2022competition, jimenez2023swe, jain2024livecodebench}.
Recent Large Language Models have demonstrated remarkable programming capabilities \citep{ge2025surveyvibecodinglarge}. Foundation model families such as GPT \citep{achiam2023gpt}, Gemini \citep{team2023gemini}, Claude \citep{anthropic2024claude3}, DeepSeek \citep{liu2025deepseek}, Kimi \citep{team2025kimi}, and Qwen \citep{yang2025qwen3} have shown impressive proficiency in generating code from natural language descriptions. However, synthesizing correct programs remains inherently challenging, requiring deep understanding of natural language processing, algorithmic design, data structure selection, and problem-solving strategies. These challenges are compounded for competitive programming and advanced software engineering tasks, where adherence to strict constraints and passing comprehensive test suites are paramount \citep{chen2021evaluating, hendrycks2021measuring, li2022competition, jimenez2023swe, jain2024livecodebench}.

% 이러한 어려움을 극복하기 위해 다양한 기법들이 발전해왔습니다. 초기 방법들은 직접 프롬프팅, Chain-of-Thought 추론, 검색 증강 생성, 다양한 인컨텍스트 예제 기반 전략을 활용했습니다. 최근 패러다임은 계획 기반 분해, 트리 탐색, 유추 추론, 다양한 에이전트 기반 아키텍처로 전환되었습니다. 이 중 \textbf{MapCoder} \citep{islam2024mapcoder}는 관심사를 분리하는 다중 에이전트 파이프라인을 도입하여 상당한 진전을 이루었습니다. 한 에이전트는 유사한 해결된 문제를 검색하고, 다른 에이전트는 실행 계획을 구성하며, 세 번째 에이전트는 코드를 생성하고, 네 번째 에이전트는 반복적인 개선을 통해 디버깅을 처리합니다. 이러한 기반을 바탕으로 \textbf{CodeSIM} \citep{islam2025codesim}은 \textit{Mental Simulation}을 두 핵심 단계에서 활용하여 최첨단 결과를 달성했습니다: (1) 계획 에이전트에서 코드 생성 전에 계획을 검증하기 위해 실행을 시뮬레이션하고, (2) 디버깅 에이전트에서 테스트 실패 시 버그를 추적하기 위해 시뮬레이션을 다시 활용합니다.
To overcome these challenges, a rich line of research has emerged. Early methods employed direct prompting \citep{chen2021evaluating}, Chain-of-Thought reasoning \citep{wei2022chain}, retrieval-augmented generation \citep{parvez2021retrieval}, and various in-context exemplar-based strategies \citep{zhang2022automatic, shum2023automatic}. Recent paradigms have shifted toward plan-based decomposition \citep{jiang2023self}, tree-searching approaches \citep{zhou2023language}, analogical reasoning \citep{yasunaga2023large}, and diverse agent-based architectures \citep{shinn2023reflexion, qian2024chatdev}. Among these, MapCoder \citep{islam2024mapcoder} marked a significant advance by introducing a multi-agent pipeline that separates concerns: one agent retrieves analogous solved problems, another constructs execution plans, a third generates code, and a fourth handles debugging through iterative refinement. Building on this foundation, CodeSIM \citep{islam2025codesim} achieved state-of-the-art results by leveraging \textit{Mental Simulation} at two critical phases: (1) in the planning agent, the LLM simulates execution to verify the plan before code generation, and (2) in the debugging agent, when code fails tests, simulation is employed again to trace through the logic and identify bugs.

% 그러나 이 접근 방식에는 근본적인 결함이 있습니다. 바로 LLM이 환각을 겪는다는 것입니다. 복잡한 알고리즘 시나리오에서 모델은 실제 프로그램 동작을 반영하지 않는 실행 추적을 ``상상''하는 경우가 많습니다. 즉, 결함이 있는 코드가 제대로 작동한다고 스스로 확신하는 희망적 사고에 빠지는 것입니다. 이는 눈을 감고 체스를 두면서 체스판의 말을 움직이지도 않고 승리를 선언하는 것과 같습니다. 우리는 이 현상을 \textbf{"Mental-Reality Gap"}이라 명명합니다: 이는 에이전트가 프로그램의 동작을 ``믿는'' 방식과 런타임에서 실제로 동작하는 방식 사이의 단절을 의미합니다. 이 간극으로 인해 에이전트는 버그가 있는 코드를 자신 있게 검증하고, 잘못된 해결책을 최종 답으로 제출하게 됩니다. 더 심각한 것은 이러한 자기 확신 검증 루프가 오류를 수정하기보다는 오히려 강화할 수 있다는 것입니다.
However, this approach has a fundamental flaw: \textbf{LLMs hallucinate}. In complex algorithmic scenarios, models frequently ``imagine'' execution traces that do not reflect actual program behavior---effectively engaging in wishful thinking where they convince themselves that flawed code works correctly \citep{huang2025survey}. This is akin to playing blindfold chess and declaring victory without ever moving the pieces on the board. We term this phenomenon the \textbf{``Mental-Reality Gap''}---the disconnect between how an agent \textit{believes} a program behaves and how it \textit{actually} behaves at runtime. This gap causes agents to confidently validate buggy code, submitting incorrect solutions as final answers. Worse, such a self-affirming verification loop can reinforce errors rather than correcting them.

% 정신-현실 간극을 메우는 가장 직접적인 방법은 코드의 동작을 상상하는 대신 실제로 코드를 실행하는 것입니다. 역사적으로 이러한 접근 방식을 가로막는 두 가지 장벽이 있었습니다. 첫 번째는 기술적 타당성입니다. LLM이 순수한 텍스트 생성기였을 때는 외부 환경과의 상호작용이 부자연스러웠고, 실행 결과를 추론에 통합하는 것이 어려웠습니다. 하지만 도구 사용과 함수 호출이 핵심 기능으로 부상하면서 이러한 장벽은 크게 해소되었습니다. 현대 에이전트는 코드를 실행하고, 오류 메시지를 해석하고, 경험적 증거를 기반으로 접근 방식을 개선할 수 있습니다 \citep{yao2022react, schick2024toolformer}. 두 번째 장벽은 더 근본적인 문제인 누락된 오라클 문제입니다. 코드를 실행하려면 예상 출력을 가진 테스트 케이스가 필요하지만, 새로운 알고리즘 문제의 경우 정확한 출력은 바로 우리가 찾고자 하는 것입니다.
The most direct way to bridge the Mental-Reality Gap is to \textit{actually execute} the code rather than imagining its behavior. Two barriers have historically prevented this approach. The first is \textbf{technical feasibility}: when LLMs were pure text generators, interaction with external environments felt unnatural, and incorporating execution results into reasoning was awkward. This barrier has largely fallen---with tool use and function calling emerging as core capabilities \citep{yao2022react, schick2024toolformer}, modern agents can execute code, interpret error messages, and refine their approach based on empirical evidence. The second barrier is more fundamental: the \textbf{Missing Oracle problem}. Executing code requires test cases with expected outputs, but for novel algorithmic problems, the correct output is precisely what we are trying to discover.

% CodeSIM의 저자들은 이러한 오라클 과제에 직접 맞섰습니다. 그들은 자기 일관성을 통해 테스트 케이스를 증강하는 방법을 모색했습니다. 즉, 테스트를 두 번 생성하고 두 응답 모두에 나타나는 테스트 케이스만 유지하는 것입니다. 그 결과 성능이 9.3%p 감소했으며, "추가 I/O 생성은 최종 결과에 부정적인 영향을 미칠 수 있는 간단한 작업이 아니다"라는 결론을 내렸습니다. 그러나 저희는 이 실패가 테스트 생성 자체에서 비롯된 것이 아니라, 실제 결과 없이 정확한 출력을 예측하려는 시도에서 비롯되었다고 주장합니다. 핵심적인 통찰력은 검증에 정확한 답이 반드시 필요하지 않다는 것입니다. 값이 아닌 속성을 확인함으로써(예: "출력 길이가 입력 길이와 같음", "결과는 입력의 순열") 오라클 없이도 정확성을 판단할 수 있습니다. 이러한 속성 기반 접근 방식은 CodeSIM이 열어둔 검증 루프를 닫습니다.
CodeSIM's authors confronted this oracle challenge directly. They explored augmenting test cases through self-consistency---generating tests twice and retaining only those appearing in both responses. The result was a \textbf{9.3 percentage point performance decline}, leading them to conclude that ``generating additional I/Os is a non-trivial task that may negatively impact final outcomes'' \citep{islam2025codesim}. However, we argue that the failure stemmed not from test generation itself, but from attempting to predict \textit{exact outputs} without ground truth. Our key insight is that verification need not require exact answers: by checking \textit{properties} rather than values (e.g., ``output length equals input length,'' ``result is a permutation of input''), we can judge correctness without an oracle. This property-based approach closes the verification loop that CodeSIM left open.

% 이러한 통찰력을 바탕으로, 우리는 정신적 추측을 실질적인 실행 피드백으로 체계적으로 대체하는 다중 에이전트 프레임워크인 SolidCoder를 제안합니다. SolidCoder의 핵심은 다섯 가지 시너지 효과를 제공하는 구성 요소를 통해 정신적-현실 간 격차(Mental-Reality Gap)를 해결하는 S.O.L.I.D. 아키텍처입니다. S.O.L.I.D. 아키텍처는 다음과 같은 다섯 가지 시너지 효과를 제공합니다. S.O.L.I.D.는 설계 시점에 버그로 이어지기 전에 경계 사례를 식별합니다. Oracle 기반 Assertions는 정확한 출력이 아닌 속성을 확인하여 실제 결과 없이도 검증을 가능하게 합니다. Live Execution은 샌드박스 코드 실행을 통해 상상을 런타임 사실로 대체합니다. I.Intermediate Simulation은 계획-코드 변환 오류를 포착합니다. D. Defensive Accumulation은 발견된 모든 실패 사례를 누적하여 회귀를 방지합니다. 우리의 기여는 다음과 같습니다.
Building on this insight, we propose \textbf{SolidCoder}, a multi-agent framework that systematically replaces mental speculation with tangible execution feedback. At its core is the \textbf{S.O.L.I.D.} architecture, which addresses the Mental-Reality Gap through five synergistic components: \textbf{S}hift-left Planning identifies edge cases at design time before they become bugs; \textbf{O}racle-based Assertions enable verification without ground truth by checking properties rather than exact outputs; \textbf{L}ive Execution replaces imagination with runtime facts through sandboxed code execution; \textbf{I}ntermediate Simulation catches plan-to-code translation errors; and \textbf{D}efensive Accumulation prevents regression by accumulating all discovered failure cases. Our contributions are:

% \begin{itemize}
% \item 기존 다중 에이전트 코드 생성 프레임워크에서 검증을 위해 정신 시뮬레이션에 의존하는 근본적인 한계인 \textbf{정신-현실 간극(Mental-Reality Gap)}을 파악하고 공식적으로 정의합니다.
% \item 정신-현실 간극을 식별하고 형식적으로 정의합니다---이는 검증을 위해 정신 시뮬레이션에 의존하는 기존 멀티-에이전트 코드 생성 프레임워크의 근본적인 한계입니다.
% \item S.O.L.I.D. 아키텍처를 통해 구체적인 실행에 검증을 기반을 두어 이러한 간극을 메우는 최초의 프레임워크인 \textbf{SolidCoder}를 제안합니다.
% \item GPT-4o를 사용하여 새로운 최첨단 pass@1 결과를 달성했습니다: HumanEval에서 \textbf{95.7\%}(+0.6\%p), CodeContests에서 \textbf{77.0\%}(+4.3\%p), APPS에서 \textbf{26.7\%}(+3.4\%p). 이러한 개선은 RL 기반 post-training을 거친 최신 모델들에서도 일관되게 나타나며, 특히 CodeContests에서 가장 큰 개선을 보입니다---정신 시뮬레이션이 실패하기 시작하지만 문제가 여전히 실행 기반 검증으로 해결 가능한 영역이기 때문입니다.
\begin{itemize}
    \item We identify and formally define the \textbf{Mental-Reality Gap}---a fundamental limitation in existing multi-agent code generation frameworks that rely on mental simulation for verification.
    \item We propose \textbf{SolidCoder}, the first framework to bridge this gap by grounding verification in concrete execution through the \textbf{S.O.L.I.D.} architecture.
    \item We achieve new state-of-the-art pass@1 results with GPT-4o: \textbf{95.7\%} on HumanEval (+0.6\%p), \textbf{77.0\%} on CodeContests (+4.3\%p), and \textbf{26.7\%} on APPS (+3.4\%p). These improvements generalize consistently across recent RL post-trained models, with the largest gains on CodeContests---where mental simulation begins to fail but problems remain tractable for execution-grounded verification.
\end{itemize}

\section{Related Work}

\begin{figure*}[t]
    \centering
    \includegraphics[width=\textwidth]{fig2.png}
    % \caption{\textbf{실제 정신-현실 간 격차} 목록 회전 문제에 대한 비교 예시. \textbf{왼쪽(CodeSIM):} 정신 시뮬레이션은 코드를 추적하여 ``출력이 예상과 일치합니다. PASS''라는 잘못된 결론을 내립니다. LLM은 버그가 있음에도 불구하고 정상적인 동작을 하는 것처럼 착각합니다. \textbf{오른쪽(SolidCoder):} 실시간 실행은 실제 코드를 실행하여 구체적인 실패를 통해 버그를 드러냅니다. \texttt{AssertionError: [3,1,2] != [2,3,1]}. 이는 현실에서의 접지 검증이 정신 시뮬레이션이 놓친 버그를 어떻게 잡아내는지 보여줍니다. (이 예시는 실시간 검증 루프에 초점을 맞춥니다. \textbf{[I]} 중간 시뮬레이션은 그림 1에 나와 있습니다.)}
    \caption{\textbf{The Mental-Reality Gap in Action.} A comparative example on a list rotation problem. \textbf{Left (CodeSIM):} Mental Simulation traces through the code and incorrectly concludes ``Output matches expected. PASS''---the LLM hallucinates correct behavior despite a bug. \textbf{Right (SolidCoder):} Live Execution runs the actual code, revealing the bug through concrete failure: \texttt{AssertionError: [3,1,2] != [2,3,1]}. This demonstrates how grounding verification in reality catches bugs that mental simulation misses. (This example focuses on the Live Verification Loop; \textbf{[I]} Intermediate Simulation is depicted in Figure~\ref{fig:1}.)}
    \label{fig:2}
\end{figure*}

% ============================================================
% [Figure 2 - Mental-Reality Gap] PPT 재작업용 최종 프롬프트
% ============================================================
%
% ========== 캔버스 설정 ==========
% 비율: 16:9 또는 4:3 슬라이드
% 배경: 흰색
% 폰트: 본문 Sans-serif (맑은고딕, Arial), 코드 Monospace (Consolas)
%
% ========== 전체 레이아웃 ==========
%
%   ┌─────────────────────────────────────────────────────────┐
%   │              [문제 박스 - 상단 중앙]                      │
%   │              "Mental-Reality Gap" 제목                   │
%   ├─────────────────────┬───────────────────────────────────┤
%   │                     │                                   │
%   │    왼쪽 컬럼         │          오른쪽 컬럼               │
%   │    (CodeSIM)        │          (SolidCoder)             │
%   │    폭: 40%          │          폭: 60%                  │
%   │                     │                                   │
%   ├─────────────────────┴───────────────────────────────────┤
%   │              [하단 인사이트 박스]                         │
%   └─────────────────────────────────────────────────────────┘
%
% ============================================================
% 섹션 1: 상단 중앙 - 문제 박스
% ============================================================
% 배경: #FFF9C4 (연한 노랑)
% 테두리: #FBC02D, 1px, 둥근모서리 8px
%
% 내용:
%   "Problem: Rotate List" (굵게, 16pt)
%   "Given a list and integer k, rotate the list right by k positions."
%   
%   Input:  [1, 2, 3], k=1
%   Output: [3, 1, 2]
%
% 아래 중앙:
%   "Mental-Reality Gap" (굵게, 28pt, 색상 #1A237E)
%
% ============================================================
% 섹션 2: 왼쪽 컬럼 - CodeSIM
% ============================================================
% 배경: #FFF8E1 (연한 크림/노랑)
% 헤더: "CodeSIM" (굵게, 20pt) + "�� Mental Simulation" (14pt)
%
% ─────────────────────────────────────
% 박스 2.1: Generated Code
% ─────────────────────────────────────
% 배경: 흰색, 테두리 #E0E0E0
%
% 헤더: "Generated Code" (굵게)
% 코드 (Monospace, 11pt):
%
%   def rotate(lst, k):
%       if not lst: return lst
%       k = k % len(lst)
%       return lst[k:] + lst[:k]   # THIS IS THE BUG
%
% ─────────────────────────────────────
% 박스 2.2: LLM Mental Trace (k=1)
% ─────────────────────────────────────
% 배경: #FFF8E1 (연한 노랑)
% 테두리: #FFE082
%
% 헤더: "�� LLM Mental Trace (k=1)" (굵게, 주황)
%
% 내용 (Monospace):
%   Input:    lst=[1,2,3], k=1
%   Step 1:   k = 1 % 3 = 1
%   Step 2:   lst[-1:] → [3]       ← LLM hallucination! (code is lst[k:])
%   Step 3:   lst[:-1] → [1, 2]    ← LLM hallucination! (code is lst[:k])
%   Step 4:   [3] + [1, 2] → [3, 1, 2]
%   Expected: [3, 1, 2]
%   Traced:   [3, 1, 2] ✅
%
% ※ "← LLM hallucination!" 부분: 이탤릭, 주황색 (#E65100)
%
% ─────────────────────────────────────
% 박스 2.3: Verdict (통과 판정)
% ─────────────────────────────────────
% 배경: #C8E6C9 (연한 초록)
%
%   ✅ PASS - Code is correct
%   Output matches expected value.
%   No further debugging needed.
%
% ─────────────────────────────────────
% 박스 2.4: HIDDEN BUG (경고)
% ─────────────────────────────────────
% 배경: #FFEBEE (연한 빨강)
% 테두리: 2px 점선 빨강 (#E53935)
%
%   ⚠️ HIDDEN BUG (not detected)
%   Test k=2: Expected [2,3,1]
%             Actual   [3,1,2] ❌
%   
%   LLM hallucinated correct behavior for k=1
%   but bug exists for other inputs.
%
% ============================================================
% 섹션 3: 오른쪽 컬럼 - SolidCoder
% ============================================================
% 배경: #E3F2FD (연한 파랑)
% 헤더: "SolidCoder" (굵게, 20pt, #1565C0) + "▶ Live Execution"
%
% ★★★ 레이아웃 (수직 흐름 권장) ★★★
%
%   ┌─────────────────┐
%   │ Generated Code  │
%   └────────┬────────┘
%            │
%   ┌────────┴────────┐
%   │ [S] Edge Cases  │ ← 작은 뱃지 스타일
%   └────────┬────────┘
%            │ k=2
%            ▼
%   ┌─────────────────┐
%   │ [O] Oracle Test │
%   └────────┬────────┘
%            │
%            ▼
%   ┌─────────────────┐
%   │ [L] Live Exec   │
%   └────────┬────────┘
%            │
%            ▼
%   ┌─────────────────┐
%   │ ❌ Bug Detected │
%   └────────┬────────┘
%            │
%            ▼
%   ┌─────────────────┐
%   │ [D] Accumulated │
%   └────────┬────────┘
%            │
%            ▼
%   ┌─────────────────┐
%   │ ✅ Fixed Code   │
%   └─────────────────┘
%
% ─────────────────────────────────────
% 박스 3.1: Generated Code (Same)
% ─────────────────────────────────────
% 배경: 흰색
%
%   def rotate(lst, k):
%       if not lst: return lst
%       k = k % len(lst)
%       return lst[k:] + lst[:k]   # THIS IS THE BUG
%
% ─────────────────────────────────────
% 박스 3.2: [S] Edge Cases Identified
% ─────────────────────────────────────
% 위치: Generated Code 옆 또는 아래
% 크기: 작은 뱃지 스타일
% 배경: #BBDEFB (연한 파랑)
%
%   [S] Shift-left Planning
%   • Empty list
%   • k=0, k > len(lst)
%   • k=2 (boundary) ← test
%
% ─────────────────────────────────────
% 박스 3.3: [O] Oracle-based Test
% ─────────────────────────────────────
% 배경: #BBDEFB
% 라벨: "[O]" 파란 뱃지
%
%   [O] Property-based Test Generation
%   
%   # Properties to verify:
%   # 1. len(output) == len(input)
%   # 2. set(output) == set(input)
%   
%   test_input = [1, 2, 3]
%   test_k = 2    # from [S] edge case
%   expected = [2, 3, 1]
%
% ─────────────────────────────────────
% 박스 3.4: [L] Live Execution
% ─────────────────────────────────────
% 배경: #263238 (다크 터미널)
% 라벨: "[L]" 파란 뱃지
%
%   [L] Sandbox Execution (5s timeout)
%   
%   >>> rotate([1,2,3], 2)
%   [3, 1, 2]                        ← 주황색
%   
%   >>> assert result == [2,3,1]
%   AssertionError: [3,1,2] != [2,3,1]   ← 빨간색
%   
%   ❌ TEST FAILED
%
% ─────────────────────────────────────
% 박스 3.5: Bug Detected
% ─────────────────────────────────────
% 배경: #FFCDD2 (연한 빨강)
%
%   ❌ Bug Detected!
%   Expected: [2, 3, 1]
%   Actual:   [3, 1, 2]
%   → Routing to Debugging Agent
%
% ─────────────────────────────────────
% 박스 3.6: [D] Accumulated Tests
% ─────────────────────────────────────
% 배경: #BBDEFB
% 라벨: "[D]" 파란 뱃지
%
%   [D] Test Suite Updated
%   
%   accumulated_inputs += {
%       ([1,2,3], k=2) → [2,3,1]
%   }
%
% ─────────────────────────────────────
% 박스 3.7: Fixed Code
% ─────────────────────────────────────
% 배경: #C8E6C9 (연한 초록)
%
%   ✅ Fixed Code
%   
%   def rotate(lst, k):
%       if not lst: return lst
%       n = len(lst)
%       k = k % n
%       return lst[n-k:] + lst[:n-k]   # ← Fixed!
%   
%   ✅ All tests pass (incl. accumulated)
%
% ============================================================
% 섹션 4: 하단 인사이트 박스
% ============================================================
% 배경: #E1F5FE
% 전체 폭, 중앙 정렬
%
%   "Mental Simulation can hallucinate correct traces for buggy code."
%   "Live Execution grounds verification in reality, catching bugs that imagination misses."
%
%   ※ "Mental Simulation"과 "Live Execution" 굵게
%
% ============================================================
% 섹션 5: 화살표 명세 (총 5개)
% ============================================================
%
% ★★★ 화살표 1: [S] → [O] ★★★
%   시작: [S] Edge Cases 박스 하단
%   끝:   [O] Oracle-based Test 박스 상단
%   스타일: 파란색 (#1976D2), 2px, 실선
%   라벨: "k=2" (화살표 옆에 표시)
%   방향: 아래로 (수직 또는 약간 대각선)
%
% ★★★ 화살표 2: [O] → [L] ★★★
%   시작: [O] 박스 하단 (또는 오른쪽)
%   끝:   [L] 박스 상단 (또는 왼쪽)
%   스타일: 파란색 (#1976D2), 2px, 실선
%   라벨: 없음
%
% ★★★ 화살표 3: [L] → Bug Detected ★★★
%   시작: [L] 박스 하단
%   끝:   Bug Detected 박스 상단
%   스타일: 빨간색 (#E53935), 2px, 실선
%   라벨: 없음 (또는 "FAIL")
%
% ★★★ 화살표 4: Bug Detected → [D] ★★★
%   시작: Bug Detected 박스 하단
%   끝:   [D] Accumulated Tests 박스 상단
%   스타일: 파란색 (#1976D2), 2px, 실선
%   라벨: 없음
%
% ★★★ 화살표 5: [D] → Fixed Code ★★★
%   시작: [D] 박스 하단 (또는 오른쪽)
%   끝:   Fixed Code 박스 상단 (또는 왼쪽)
%   스타일: 초록색 (#43A047), 2px, 실선
%   라벨: 없음
%
% ============================================================
% 섹션 6: 색상 팔레트
% ============================================================
%
% 배경색:
%   CodeSIM 영역:    #FFF8E1 (연한 크림)
%   SolidCoder 영역: #E3F2FD (연한 파랑)
%   문제 박스:       #FFF9C4 (연한 노랑)
%   인사이트 박스:   #E1F5FE (연한 하늘)
%
% 컴포넌트 박스:
%   S.O.L.I.D. 박스: #BBDEFB (연한 파랑)
%   터미널 [L]:      #263238 (다크)
%   성공 박스:       #C8E6C9 (연한 초록)
%   실패/버그 박스:  #FFCDD2 (연한 빨강)
%   경고 박스:       #FFEBEE (더 연한 빨강)
%
% 텍스트:
%   S.O.L.I.D. 라벨: #1565C0 (파란 굵게)
%   성공 텍스트:     #2E7D32 (초록)
%   실패 텍스트:     #C62828 (빨강)
%   환각 주석:       #E65100 (주황 이탤릭)
%
% 화살표:
%   일반 흐름:       #1976D2 (파랑)
%   실패 흐름:       #E53935 (빨강)
%   성공 흐름:       #43A047 (초록)
%
% ============================================================
% 섹션 7: 최종 체크리스트
% ============================================================
%
% 코드 로직:
%   ☐ Generated Code: lst[k:] + lst[:k] (버그 - 좌회전)
%   ☐ Mental Trace: lst[-1:], lst[:-1] + 환각 주석
%   ☐ Live Execution: rotate([1,2,3], 2) → [3,1,2]
%   ☐ AssertionError: [3,1,2] != [2,3,1]
%   ☐ Fixed Code: lst[n-k:] + lst[:n-k]
%
% S.O.L.I.D. 라벨:
%   ☐ [S] Shift-left Planning
%   ☐ [O] Oracle-based Test (NOT [S]!)
%   ☐ [L] Live Execution + "5s timeout"
%   ☐ [D] "Test Suite Updated" (NOT "Suited"!)
%
% 화살표 (5개):
%   ☐ [S] → [O] (파랑, k=2 라벨)
%   ☐ [O] → [L] (파랑)
%   ☐ [L] → Bug Detected (빨강)
%   ☐ Bug Detected → [D] (파랑)
%   ☐ [D] → Fixed Code (초록)
%
% 기타:
%   ☐ CodeSIM 배경 크림/노랑 계열
%   ☐ SolidCoder 배경 파랑 계열
%   ☐ 하단 인사이트 문구 포함
%   ☐ "# THIS IS THE BUG" 주석 표시
%   ☐ "# ← Fixed!" 주석 표시
%
% ============================================================

\subsection{Program Synthesis and LLMs for Code}
% 프로그램 합성은 AI 분야에서 오랜 역사를 가지고 있으며, 초기 검색 기반 방법(manna1971toward, gulwani2011automating)에서 신경망 기반 방법(feng2020codebert, ahmad2021unified)으로 발전해 왔습니다. 대규모 언어 모델(LLM)의 등장은 이러한 환경을 근본적으로 변화시켰습니다.
Program synthesis has a long-standing history in AI, evolving from early search-based methods \citep{manna1971toward, gulwani2011automating} to neural approaches \citep{feng2020codebert, ahmad2021unified}. The advent of Large Language Models (LLMs) has fundamentally transformed this landscape.

% Codex \citep{chen2021evaluating}는 코드용 LLM을 개척했습니다. GPT-4 \citep{achiam2023gpt}와 GPT-4o \citep{hurst2024gpt}는 실제 작업에 적합한 고품질 코드를 생성하여 전환점을 마련했습니다. Gemini \citep{team2023gemini}와 Claude \citep{anthropic2024claude3}를 포함한 다른 모델들도 유사한 성능을 보여주었습니다. 오픈 소스 커뮤니티는 빠르게 격차를 좁혀왔습니다. Llama 계열 \citep{roziere2023code, dubey2024llama}, DeepSeek-Coder \citep{guo2024deepseek} 및 DeepSeek-V3 \citep{liu2024deepseek}, 그리고 Qwen-2.5 Coder \citep{hui2024qwen2} 및 Qwen-3 \citep{yang2025qwen3}는 특화된 코드 모델이 상용 모델과 경쟁할 수 있음을 보여주었습니다. 최근에는 OpenAI o1 \citep{jaech2024openai}을 시작으로 DeepSeek-R1 \citep{guo2025deepseek}에 이르기까지 추론 모델들이 실행 피드백을 훈련 신호로 사용하는 검증 가능한 보상 기반 강화 학습(RLVR) \citep{wang2025reinforcementlearningreasoninglarge}과 같은 기술을 통해 확장된 사고 연쇄 추론을 통해 코드 생성을 더욱 발전시켰습니다. 이러한 모델들의 코딩 및 명령 추종 능력이 지속적으로 향상됨에 따라, 단순한 코드 생성기를 넘어 자율 에이전트의 "두뇌" 역할을 점점 더 수행하고 있습니다.
Codex \citep{chen2021evaluating} pioneered LLMs for code. GPT-4 \citep{achiam2023gpt} and GPT-4o \citep{hurst2024gpt} marked a turning point by generating production-quality code suitable for real-world tasks. Other models including Gemini \citep{team2023gemini} and Claude \citep{anthropic2024claude3} have demonstrated comparable capabilities. The open-source community has rapidly closed the gap: the Llama family \citep{roziere2023code, dubey2024llama}, DeepSeek-Coder \citep{guo2024deepseek} and DeepSeek-V3 \citep{liu2024deepseek}, and the Qwen-2.5 Coder \citep{hui2024qwen2} \& Qwen-3 \citep{yang2025qwen3} showed that specialized code models can rival proprietary alternatives. More recently, reasoning models---starting with OpenAI o1 \citep{jaech2024openai} and followed by DeepSeek-R1 \citep{guo2025deepseek}---have further advanced code generation through extended chain-of-thought reasoning, with techniques like Reinforcement Learning from Verifiable Rewards (RLVR) \citep{wang2025reinforcementlearningreasoninglarge} using execution feedback as a training signal. As these models' coding and instruction-following capabilities continue to grow, they are increasingly serving as the ``brain'' for autonomous agents rather than mere code generators.

\subsection{Approaches to LLM Code Generation}
% 유능한 LLM을 기반으로, 이제 문제는 LLM을 어떻게 최적으로 조율할 것인가입니다. 코드 생성을 향상시키기 위해 다양한 프롬프트 전략이 개발되었습니다. 사고 연쇄(CoT)는 복잡한 문제를 중간 추론 단계로 분해합니다. 검색 증강 생성(Retrieval-augmented generation)은 유사한 사례를 활용하여 생성을 안내합니다. 자가 계획(Self planning)은 코딩 전에 알고리즘 계획을 생성합니다. 유추 추론(Analogical reasoning)은 관련 문제에서 해결책을 회상하고 수정합니다. ReAct는 추론과 액션 실행을 병행합니다. 이러한 기법들은 단일 패스 생성을 개선하지만, LLM이 실제 실행을 통해 자체 출력을 검증할 수 없다는 한계를 가지고 있습니다.
With capable LLMs as the foundation, the question becomes \textit{how} to best orchestrate them. Various prompting strategies have been developed to enhance code generation. Chain-of-Thought (CoT) \citep{wei2022chain} decomposes complex problems into intermediate reasoning steps. Retrieval-augmented generation \citep{parvez2021retrieval} leverages similar examples to guide generation. Self-planning \citep{jiang2023self} generates algorithmic plans before coding. Analogical reasoning \citep{yasunaga2023large} recalls and adapts solutions from related problems. ReAct \citep{yao2022react} interleaves reasoning with action execution. While these techniques improve single-pass generation, they remain limited by the LLM's inability to verify its own outputs through actual execution.

% 이러한 한계를 해결하기 위해 최근 연구는 소프트웨어 개발 수명 주기를 시뮬레이션하는 다중 에이전트 프레임워크로 전환되었습니다. MapCoder는 검색, 계획, 코딩 및 디버깅 에이전트로 구성된 파이프라인을 도입하여 단일 에이전트 기준선보다 훨씬 우수한 성능을 보였습니다. CodeSIM은 시뮬레이션 기반 계획 및 디버깅을 통합하여 이를 더욱 향상시켰으며, 에이전트가 단계별 실행 추적을 통해 로직을 검증할 수 있도록 했습니다. Reflexion과 LATS는 각각 반복적인 자체 수정 및 트리 탐색을 활용합니다. 그러나 이러한 프레임워크는 여전히 LLM의 내부 추론에 의존하여 검증하기 때문에 실행 추적 오류에 취약합니다.
To address this limitation, recent research has shifted towards multi-agent frameworks that simulate the software development lifecycle. MapCoder \citep{islam2024mapcoder} introduced a pipeline of retrieval, planning, coding, and debugging agents, significantly outperforming single-agent baselines. CodeSIM \citep{islam2025codesim} further enhanced this by incorporating simulation-driven planning and debugging, allowing agents to verify their logic through step-by-step execution traces. Other approaches like Reflexion \citep{shinn2023reflexion} and LATS \citep{zhou2023language} leverage iterative self-correction and tree search, respectively. However, these frameworks still rely on the LLM's internal reasoning for verification, making them susceptible to hallucinated execution traces.

% 별도의 작업 라인에서 실제 코드 실행을 디버깅에 활용하는 방안을 모색했습니다. LDB \citep{zhong2024debuglikehumanlarge}는 프로그램을 기본 블록으로 분할하고 런타임 중 중간 변수 값을 추적하여 불일치를 단계별로 식별합니다. MGDebugger \citep{shi2025codecorrectnessclosingmile}은 계층적 접근 방식을 사용하여 코드를 하위 함수 트리로 분해하고 상향식 디버깅을 수행합니다. 이러한 실행 기반 디버거는 실제로는 검증 기반이지만, 초기 코드 생성 후 두 번째 단계(\textit{second pass})로 작동하며 오류 감지를 위해 실제 검증 테스트 케이스가 필요합니다.
A separate line of work has explored leveraging actual code execution for debugging. LDB \citep{zhong2024debuglikehumanlarge} segments programs into basic blocks and tracks intermediate variable values during runtime, identifying discrepancies step-by-step. MGDebugger \citep{shi2025codecorrectnessclosingmile} takes a hierarchical approach, decomposing code into a tree of sub-functions and debugging bottom-up. While these execution-based debuggers ground verification in reality, they operate as a \textit{second pass} after initial code generation and require ground-truth test cases to detect errors.

% 위에서 논의된 접근 방식은 두 가지 범주로 나뉩니다. (1) MapCoder 및 CodeSIM과 같은 다중 에이전트 프레임워크는 코드 생성에는 탁월하지만 검증을 위해 LLM의 내부 추론에 의존합니다. (2) LDB 및 MGDebugger와 같은 외부 디버거는 실제 실행을 활용하지만 생성 루프 외부의 사후 단계처럼 작동합니다. 두 가지 모두 심각한 간극을 남깁니다. 다중 에이전트 프레임워크는 환각적인 검증에 취약하고, 외부 디버거는 실제 테스트 케이스를 요구하며 버그 발생 자체를 방지할 수 없습니다.
The approaches discussed above fall into two categories: (1) multi-agent frameworks like MapCoder and CodeSIM excel at code generation but rely on the LLM's internal reasoning for verification, and (2) external debuggers like LDB and MGDebugger leverage actual execution but operate as post-hoc steps outside the generation loop. Both leave critical gaps---multi-agent frameworks are prone to hallucinated verification, while external debuggers require ground-truth test cases and cannot prevent bugs from forming in the first place.

\section{SolidCoder Methodology}
\label{sec:methodology}

% SolidCoder는 두 패러다임의 장점을 통합하는 동시에 그 한계를 해결합니다. 다중 에이전트 프레임워크처럼 통합된 생성 파이프라인 내에서 작동하며, 외부 디버거처럼 실제 실행을 활용합니다. CodeSIM의 세 에이전트 아키텍처(계획, 코딩, 디버깅)를 기반으로 하며, S.O.L.I.D. 프레임워크를 통해 정신-현실 간 격차를 체계적으로 해결합니다. 핵심 통찰은 간단합니다: 상상하지 말고 실행하세요.
SolidCoder unifies the strengths of both paradigms while addressing their limitations. Like multi-agent frameworks, we operate within an integrated generation pipeline; like external debuggers, we leverage actual execution. Building upon the three-agent architecture of CodeSIM---comprising Planning, Coding, and Debugging agents---we introduce the \textbf{S.O.L.I.D.} framework to systematically address the Mental-Reality Gap. Our guiding principle is straightforward: \textit{don't imagine---execute}. Rather than relying on mental simulation for verification, we integrate concrete execution feedback directly into the generation loop.

% SolidCoder의 주요 혁신은 세 가지입니다. (1) 사후 적용 대신 실시간 실행을 생성 루프에 통합합니다. (2) 오라클 기반 어설션을 통해 실제 정답 없이 속성 검증을 가능하게 합니다. (3) 방어적 누적 기능을 통해 발견된 버그가 재발하지 않도록 보장합니다. 표 1은 SolidCoder를 기존 코드 생성 방식과 비교합니다.
Our key innovations are threefold: (1) we integrate live execution \textit{into} the generation loop rather than applying it post-hoc, (2) oracle-based assertions enable property verification without ground-truth answers, and (3) Defensive Accumulation ensures discovered bugs never regress. Table~\ref{tab:comparison} compares SolidCoder with existing code generation approaches.

\begin{table}[t]
    \centering
    \resizebox{\columnwidth}{!}{%
    \begin{tabular}{@{}lccccccc@{}}
        \toprule
        & \multicolumn{4}{c}{\textit{Generation}} & \multicolumn{3}{c}{\textit{Verification}} \\
        \cmidrule(lr){2-5} \cmidrule(l){6-8}
        \textbf{Approach} &
        \makecell{Exemplar\\Retrieval} &
        Plan &
        \makecell{Edge\\Aware} &
        \makecell{Mental\\Simulation} &
        \makecell{Live\\Execution} &
        Debug &
        \makecell{Defensive\\Accumulation} \\
        \midrule
        Reflexion     & \xmark & \xmark & \xmark & \xmark & \xmark & \cmark & \xmark \\
        Self-planning & \xmark & \cmark & \xmark & \xmark & \xmark & \xmark & \xmark \\
        Analogical    & \cmark & \xmark & \xmark & \xmark & \xmark & \xmark & \xmark \\
        LATS          & \xmark & \cmark & \xmark & \xmark & \xmark & \cmark & \xmark \\
        MapCoder      & \cmark & \cmark & \xmark & \xmark & \xmark & \cmark & \xmark \\
        CodeSIM       & \cmark & \cmark & \xmark & \cmark & \xmark & \cmark & \xmark \\
        \midrule
        \rowcolor{LightBlue}
        \textbf{SolidCoder} & \cmark & \cmark & \cmark & \cmark & \cmark & \cmark & \cmark \\
        \bottomrule
    \end{tabular}}
    \caption{Comparison of code generation approaches. SolidCoder uniquely combines all capabilities.}
    \label{tab:comparison}
\end{table}

\subsection{S.O.L.I.D. Components}

% \paragraph{Shift-left Planning (S).} 기존 방식은 디버깅 중에 예외 상황을 반응적으로 처리합니다. SolidCoder는 알고리즘 계획을 수립하기 전에 잠재적 실패 사례를 식별하기 위해 Planning Agent를 확장합니다. LLM에 다음과 같이 질문합니다. \textit{``어떤 최악의 입력이 순진한 해결책을 망칠 수 있는가?''} 이러한 예외 상황(빈 입력, 경계값, 코너 사례)은 계획 프롬프트에 통합되어 처음부터 강력한 알고리즘 설계를 강제합니다.
\paragraph{Shift-left Planning (S).} Traditional approaches address edge cases reactively during debugging. SolidCoder extends the Planning Agent to identify potential failure cases \textit{before} formulating the algorithmic plan. We prompt the LLM: \textit{``What worst-case inputs could break a naive solution?''} These edge cases (empty inputs, boundary values, corner cases) are then incorporated into the planning prompt, forcing robust algorithm design from the outset.

% \paragraph{오라클 기반 단언(O).} 테스트 생성의 근본적인 과제는 새로운 입력에 대한 정확한 출력을 결정하는 ``오라클 문제''입니다. Mental Simulation은 LLM이 출력을 예측하도록 하여 환각 위험을 다시 발생시킴으로써 이 문제를 회피합니다. SolidCoder는 \textit{속성 기반 검증}을 통해 이 문제를 해결합니다. 정확한 출력을 확인하는 대신, 결과가 도메인 불변 속성(예: 정렬 함수가 길이를 유지하고, 순서를 유지하며, 순열을 생성하는지)을 충족하는지 검증합니다. 이를 통해 검증은 ``이 출력이 정확한가?''에서 ``이 출력이 필요한 속성을 충족하는가?''로 전환되며, 이는 실행을 통해 답할 수 있습니다.
\paragraph{Oracle-based Assertions (O).} A fundamental challenge in test generation is the ``oracle problem''---determining correct outputs for novel inputs. Mental Simulation sidesteps this by having the LLM predict outputs, reintroducing hallucination risk. SolidCoder addresses this through \textit{property-based verification}: rather than checking exact outputs, we verify that results satisfy domain-invariant properties (e.g., a sorting function preserves length, maintains ordering, and produces a permutation). This transforms verification from ``Is this output correct?'' to ``Does this output satisfy required properties?''---answerable through execution.

% \paragraph{실시간 실행(L).} LLM에 코드 동작을 상상하도록 요청하는 대신, LLM은 코드를 실행합니다. 멘탈 시뮬레이션은 잠재적으로 수백 개의 실행 단계에 걸쳐 정확한 상태를 유지해야 하며, 복잡성이 증가함에 따라 오류 발생 가능성이 높아집니다. 실시간 실행은 상태 추적을 실제 인터프리터에 위임하여 이러한 문제를 해결합니다. 생성된 코드는 샌드박스 환경(5초 제한 시간, 파일 시스템 접근 제한)에서 실행되며, 어설션 실패 또는 런타임 오류 발생 시 디버깅으로 결과를 전달합니다.
\paragraph{Live Execution (L).} Rather than asking the LLM to imagine code behavior, we run it. Mental Simulation requires maintaining accurate state across potentially hundreds of execution steps---increasingly error-prone as complexity grows. Live Execution eliminates this by delegating state tracking to an actual interpreter. We execute generated code in a sandboxed environment (5-second timeout, restricted filesystem access), routing results to debugging on assertion failures or runtime errors. \textit{Note:} Live Execution currently supports Python; for other languages, the system falls back to mental simulation.

% \paragraph{중간 시뮬레이션(I).} 검증된 계획이 정확한 코드를 보장하지는 않습니다. 변환 오류(하나씩 어긋나는 실수, 잘못된 연산자 우선순위)가 발생할 수 있습니다. 코드 생성 직후, 그리고 실시간 실행 전에 LLM이 샘플 입력에 대해 코드를 단계별로 추적하도록 합니다. 이는 전략적으로 Mental Simulation을 보완 필터로 \textit{유지합니다}. 추적된 출력이 예상 출력과 모순되는 경우, 실행 오버헤드가 발생하기 전에 버그의 증거를 확보할 수 있습니다.
\paragraph{Intermediate Simulation (I).} A verified plan does not guarantee correct code---translation errors (off-by-one mistakes, incorrect operator precedence) can occur. Immediately after code generation and before Live Execution, we prompt the LLM to trace the code step-by-step on a sample input. This strategically \textit{retains} Mental Simulation as a complementary filter: if the traced output contradicts the expected output, we have evidence of a bug before incurring execution overhead.

% \paragraph{방어적 누적(D).} 반복적 디버깅은 회귀 위험을 야기합니다. 버그 A를 수정하면 이전 반복에서 발생한 버그 B가 다시 발생할 수 있습니다. 방어적 누적은 디버깅 과정 전반에 걸쳐 확장되는 지속적인 테스트 모음을 유지합니다. 라이브 실행에서 실패한 테스트 케이스를 발견할 때마다 해당 입력-어설션 쌍은 \texttt{누적된\_테스트} 컬렉션에 추가됩니다. 중요한 것은 모든 누적 테스트가 코드 수정 후 재실행되어 단조성을 보장한다는 것입니다. 코드가 테스트를 통과하면 회귀가 즉시 감지됩니다.
\paragraph{Defensive Accumulation (D).} Iterative debugging introduces regression risk: fixing bug A may reintroduce bug B from a previous iteration. Defensive Accumulation maintains a persistent test suite that grows throughout debugging. Whenever Live Execution discovers a failing test case, that input-assertion pair joins an \texttt{accumulated\_tests} collection. Crucially, \textit{all} accumulated tests are re-executed after every code modification, providing a monotonicity guarantee: once code passes a test, regression triggers immediate detection.

\subsection{Integrated Workflow}
% 다섯 가지 구성 요소는 Appendix~\ref{app:algorithm}의 Algorithm~\ref{alg:solidcoder}에 나와 있는 워크플로에 통합됩니다. Planning Agent는 Plan을 생성하기 전에 \textbf{S}를 적용하여 경계 조건을 식별합니다. Coding Agent는 Plan을 코드로 변환한 후 \textbf{I}를 통해 변환 오류를 포착합니다. 코드는 Live Verification 루프에 진입합니다. \textbf{O}는 속성 기반 테스트를 생성하고, \textbf{L}은 Ground-Truth 피드백을 제공하며, \textbf{D}는 실패한 테스트를 누적합니다. 알고리즘은 코드가 모든 샘플 I/O와 누적된 테스트를 통과하거나, (CodeSIM처럼) $p=5$개의 Planning과 $d=5$개의 디버깅 반복을 완료하면 종료됩니다.
The five components integrate into the workflow detailed in Algorithm~\ref{alg:solidcoder} in Appendix~\ref{app:algorithm}. The Planning Agent applies \textbf{S} to identify edge cases before generating a plan. The Coding Agent translates the plan to code, followed by \textbf{I} to catch translation errors. The code enters the Live Verification loop: \textbf{O} generates property-based tests, \textbf{L} provides ground-truth feedback, and \textbf{D} accumulates failing tests. The algorithm terminates when code passes all sample I/Os and accumulated tests, or after exhausting $p=5$ planning and $d=5$ debugging iterations (following CodeSIM).

\section{Experiments}

\subsection{Datasets}
% 다양한 난이도를 나타내는 세 가지 벤치마크를 선택했습니다.
    % \begin{itemize}[leftmargin=*, itemsep=2pt, topsep=2pt]
    % \item \textbf{HumanEval} \citep{chen2021evaluating}: 짧은 docstring 사양을 가진 함수 수준 문제 164개(대개 1~2줄로 해결 가능) (\textit{easy}).
    % \item \textbf{CodeContests} \citep{li2022competition}: 정확한 I/O 사양과 문제당 약 200개의 숨겨진 테스트를 포함하는 알고리즘 추론이 필요한 Codeforces 문제 165개(\textit{medium}).
    % \item \textbf{APPS} \citep{hendrycks2021measuring}: 모호한 자연어 설명을 가진 복잡한 규칙 기반 작업을 특징으로 하는 문제 150개(\textit{hard}).
% \end{itemize}
We select three benchmarks representing different difficulty levels:
\begin{itemize}[leftmargin=*, itemsep=2pt, topsep=2pt]
    \item \textbf{HumanEval} \citep{chen2021evaluating}: 164 function-level problems with short docstring specifications, often solvable in 1--2 lines (\textit{easy}).
    \item \textbf{CodeContests} \citep{li2022competition}: 165 Codeforces problems requiring algorithmic reasoning with precise I/O specifications and ${\sim}$200 hidden tests per problem (\textit{medium}).
    \item \textbf{APPS} \citep{hendrycks2021measuring}: 150 problems featuring complex rule-based tasks with ambiguous natural language descriptions (\textit{hard}).
\end{itemize}

\subsection{Baselines and Metric}
% CodeSIM의 기준 비교군을 그대로 사용하되, Reflexion/LATS는 CodeSIM 코드베이스에서 동일 조건으로 재현 가능한 데이터/코드가 공개되지 않아 공정성을 위해 제외했습니다. 포함된 방법은 Direct, CoT, Self-Planning, yasunaga2023self, 검색 증강(analogical), 그리고 다중 에이전트(MapCoder, CodeSIM)이며, 평가는 pass@1으로 첫 시도가 모든 숨겨진 테스트를 통과하면 성공으로 간주합니다.
We adopt the comparison set from CodeSIM \citep{islam2025codesim}, excluding Reflexion \citep{shinn2023reflexion} and LATS \citep{zhou2023language} because they are not reproducible under the publicly released CodeSIM codebase (no matching code/data for identical conditions). Included baselines are prompting methods (Direct, CoT \citep{wei2022chain}, Self-Planning \citep{jiang2023self}), retrieval-augmented approaches (Analogical \citep{yasunaga2023large}), and multi-agent frameworks: MapCoder \citep{islam2024mapcoder} and CodeSIM \citep{islam2025codesim}. We report \textbf{pass@1}: a problem is solved if the first generated solution passes all hidden test cases.

\subsection{Experimental Setup}
\label{sec:experimental_setup}
% 세 가지 모델에서 평가한다: (1) GPT-4o—기존 연구와의 직접 비교를 위한 primary baseline, (2) GPT-OSS-120B—RLVR로 post-training된 GPT 아키텍처 기반 오픈소스 모델, (3) Grok-4.1-Fast—RL post-training을 거친 non-GPT 프론티어 모델로 수학/코드에서 강한 성능. 이 선택은 proprietary/open-source, pre-/post-RLVR, GPT/non-GPT의 세 축에서 일반화를 검증한다. 
We evaluate on three models: \textbf{GPT-4o} (\texttt{gpt-4o-2024-08-06}) as our primary baseline for direct comparison with prior work \citep{islam2025codesim, islam2024mapcoder}; \textbf{GPT-OSS-120B} \citep{agarwal2025gpt}, an open-source GPT model post-trained with RLVR \citep{wang2025reinforcementlearningreasoninglarge}; and \textbf{Grok-4.1-Fast} \citep{xai2025grok4card}, a non-GPT frontier model also post-trained with RL, with strong math and coding performance. This selection tests generalization across three axes: proprietary vs.\ open-source, pre- vs.\ post-RLVR, and GPT vs.\ non-GPT architectures.

% GPT-4o는 temperature τ=0, top-p=0.95를 사용하며, HumanEval 베이스라인은 CodeSIM 논문에서, CodeContests/APPS는 공식 코드베이스로 재현한다. GPT-OSS-120B와 Grok-4.1-Fast는 OpenRouter API의 기본 설정을 사용한다. 하이퍼파라미터는 CodeSIM을 따른다 (p=5, d=5); SolidCoder assumption breaking에는 a=3을 사용한다.
For GPT-4o, we use temperature $\tau{=}0$ and top-$p{=}0.95$, adopting HumanEval baselines from the CodeSIM paper and reproducing CodeContests/APPS with the official codebase.\footnote{\url{https://github.com/kagnlp/CodeGenerator}} The additional models use default configurations from their OpenRouter\footnote{\url{https://openrouter.ai}} API pages. All experiments follow CodeSIM hyperparameters ($p{=}5$ planning iterations, $d{=}5$ debugging iterations); SolidCoder uses $a{=}3$ for assumption breaking.

\section{Results}

\subsection{Main Results}

\begin{table*}[t]
    \centering
    \small
    \renewcommand{\arraystretch}{0.92}
    \setlength{\tabcolsep}{4pt}
    \resizebox{\textwidth}{!}{%
    \begin{tabular}{@{}l|ccccc|cc@{}}
        \toprule
        \textbf{Model} & \textbf{Direct} & \textbf{CoT} & \textbf{SelfPlanning} & \textbf{Analogical} & \textbf{MapCoder} & \textbf{CodeSIM} & \textbf{SolidCoder} \\
        \midrule
        \rowcolor{green!15}
        \multicolumn{8}{c}{\textbf{HumanEval}~\citep{chen2021evaluating} {\small\textcolor{green!60!black}{(EASY)}}} \\
        GPT-4o$^*$ & 90.2 & 90.9 & 89.0 & 88.4 & 90.2 & 95.1 & \textbf{95.7} \\
        GPT-OSS-120B & 69.5 & 96.3 & 90.8 & 89.0 & 61.0 & \textbf{98.2} & \textbf{98.2} \\
        Grok-4.1-Fast & 88.4 & 96.9 & 96.9 & 95.7 & 95.7 & \textbf{97.6} & \textbf{97.6} \\
        \cmidrule{1-8}
        \textit{Average} & \textit{82.7} & \textit{94.7} & \textit{92.2} & \textit{91.0} & \textit{82.3} & \textit{97.0} & \textit{\textbf{97.2}} \\
        \midrule
        \rowcolor{yellow!20}
        \multicolumn{8}{c}{\textbf{CodeContests}~\citep{li2022competition} {\small\textcolor{orange!80!black}{(MEDIUM)}}} \\
        GPT-4o$^*$ & 42.4 & 44.2 & 49.1 & 30.3 & 69.1 & 72.7 & \textbf{77.0} \\
        GPT-OSS-120B & 75.8 & 75.2 & 75.2 & 75.2 & 44.8 & 87.9 & \textbf{92.1} \\
        Grok-4.1-Fast & 79.4 & 85.4 & 81.2 & 77.0 & 83.6 & 95.2 & \textbf{98.2} \\
        \cmidrule{1-8}
        \textit{Average} & \textit{65.9} & \textit{68.3} & \textit{68.5} & \textit{60.8} & \textit{65.8} & \textit{85.3} & \textit{\textbf{89.1}} \\
        \midrule
        \rowcolor{red!15}
        \multicolumn{8}{c}{\textbf{APPS}~\citep{hendrycks2021measuring} {\small\textcolor{red!70!black}{(HARD)}}} \\
        GPT-4o$^*$ & 10.7 & 17.3 & 14.7 & 14.0 & 20.7 & 23.3 & \textbf{26.7} \\
        GPT-OSS-120B & 35.3 & 32.7 & 34.7 & 30.7 & 24.0 & 39.3 & \textbf{40.7} \\
        Grok-4.1-Fast & 37.3 & 34.7 & 37.3 & 36.0 & 33.3 & 41.3 & \textbf{42.0} \\
        \cmidrule{1-8}
        \textit{Average} & \textit{27.8} & \textit{28.2} & \textit{28.9} & \textit{26.9} & \textit{26.0} & \textit{34.6} & \textit{\textbf{36.5}} \\
        \bottomrule
    \end{tabular}}
    % 7가지 전략과 3가지 모델에 걸쳐 Pass@1 (%)을 측정했습니다. 벤치마크는 난이도 순으로 정렬되어 있습니다. GPT-4o는 기존 연구와의 비교를 위해 사용되었으며, 나머지 모델은 강화 학습(RL)을 통해 학습되었습니다. SolidCoder는 CodeSIM보다 일관적으로 우수한 성능을 보였으며, CodeContests에서 가장 큰 성능 향상(평균 +3.8%)을 기록했습니다.
    \caption{Pass@1 (\%) across seven strategies and three models. Benchmarks are ordered by difficulty. $^*$GPT-4o enables comparison with prior work; the remaining models are RL post-trained. SolidCoder consistently outperforms CodeSIM, with the largest gain on CodeContests (+3.8\%p avg). \textbf{Bold}: best per row.}
    \label{tab:main_and_multimodel}
\end{table*}

% 표~\ref{tab:main_and_multimodel}은 세 가지 벤치마크에서 결과를 보여준다. SolidCoder는 9개 모델-벤치마크 조합 모두에서 CodeSIM과 동등하거나 더 나은 성능을 보인다. GPT-4o에서 SolidCoder는 HumanEval 95.7%(+0.6%p), CodeContests 77.0%(+4.3%p), APPS 26.7%(+3.4%p)를 달성하여 CodeSIM 논문과 직접 비교가 가능하다.
Table~\ref{tab:main_and_multimodel} presents results across three benchmarks. SolidCoder matches or outperforms CodeSIM in all nine model-benchmark combinations. On GPT-4o, SolidCoder achieves \textbf{95.7\%} on HumanEval (+0.6\%p), \textbf{77.0\%} on CodeContests (+4.3\%p), and \textbf{26.7\%} on APPS (+3.4\%p)---enabling direct comparison with the CodeSIM paper~\citep{islam2025codesim}.

% 개선폭은 벤치마크 난이도에 따라 체계적인 패턴을 보인다. HumanEval은 현재 LLM들이 97%+ 수준에 도달한 포화 상태로, 추가 검증의 효과가 제한적이다. 반면 CodeContests에서는 모든 모델이 일관되게 가장 큰 개선을 보인다: GPT-4o +4.3%p, GPT-OSS-120B +4.2%p, Grok-4.1-Fast +3.0%p (평균 +3.8%p). 이는 CodeContests가 난이도 측면에서 "sweet spot"에 해당하기 때문이다---mental simulation이 실패하기 시작할 만큼 복잡하지만, 실행 기반 검증으로 오류를 효과적으로 수정할 수 있을 만큼 tractable하다. 특히 Grok-4.1-Fast는 SolidCoder와 함께 98.2%를 달성하여 이 경쟁 프로그래밍 벤치마크에서 거의 천장에 도달했다.
The pattern of improvements reveals a systematic relationship with benchmark difficulty. HumanEval, where current LLMs already exceed 97\% average performance, shows minimal gains---the benchmark is approaching saturation, leaving little room for verification-based improvements. CodeContests, by contrast, exhibits the largest and most consistent gains across all models: GPT-4o +4.3\%p, GPT-OSS-120B +4.2\%p, and Grok-4.1-Fast +3.0\%p (averaging +3.8\%p). We attribute this to CodeContests occupying a ``sweet spot'' in difficulty---problems are complex enough that mental simulation becomes unreliable, yet remain tractable enough for execution-grounded verification to effectively catch and correct errors. Notably, Grok-4.1-Fast with SolidCoder reaches \textbf{98.2\%}, approaching ceiling performance on this competitive programming benchmark.

% APPS에서는 모델에 따라 개선폭이 다르다. GPT-4o는 +3.4%p로 상당한 개선을 보이지만, RL post-trained 모델들(GPT-OSS-120B +1.4%p, Grok-4.1-Fast +0.7%p)은 더 작은 개선을 보인다. 이는 두 가지 요인으로 설명된다. 첫째, RL post-trained 모델은 학습 과정에서 코드 실행 피드백을 이미 활용하여 일부 검증 능력을 내재화했다. 둘째, APPS의 극단적 난이도로 인해 병목이 검증에서 생성으로 이동한다---완벽한 검증도 근본적으로 결함 있는 알고리즘을 구제할 수 없다.
On APPS, improvements vary by model: GPT-4o gains +3.4\%p, while RL post-trained models show smaller improvements (GPT-OSS-120B +1.4\%p, Grok-4.1-Fast +0.7\%p). Two factors explain this divergence. First, RL post-trained models have already been optimized using execution feedback during training, potentially internalizing some verification capabilities. Second, APPS poses such extreme difficulty that the bottleneck shifts from \textit{verification} to \textit{generation}---even perfect verification cannot salvage fundamentally flawed algorithms.

% 흥미로운 점은 RLVR 모델에서도 SolidCoder가 여전히 효과적이라는 것이다. RLVR 모델은 학습 시 코드 실행 피드백을 활용하여 처음부터 더 정확한 코드를 생성하도록 학습되었으므로, 검증 단계에서 잡아야 할 오류 자체가 줄어들어 SolidCoder의 추가 효과가 감소할 것으로 예상할 수 있다. 그러나 GPT-OSS-120B와 Grok-4.1-Fast는 CodeContests에서 각각 +4.2%p, +3.0%p의 상당한 개선을 보인다. 이는 RLVR이 코드 생성 능력을 향상시키더라도, 추론 시점에서 모델이 자신의 코드를 평가할 때는 여전히 mental simulation에 의존하며, SolidCoder의 live execution이 이를 대체하여 추가적인 가치를 제공함을 보여준다. 전체적으로, 이 결과는 Mental-Reality Gap이 특정 모델의 약점이 아니라 mental simulation의 근본적 한계임을 검증하며, 실행 기반 검증은 모델 계열과 학습 패러다임에 관계없이 일관된 개선을 제공한다.
An interesting finding concerns RL post-trained models. Since RLVR models are trained with execution feedback to generate more accurate code from the outset, one might expect fewer errors to catch during verification, reducing SolidCoder's marginal benefit. However, GPT-OSS-120B and Grok-4.1-Fast still show substantial gains on CodeContests (+4.2\%p and +3.0\%p). This indicates that even when RLVR improves code generation, models still rely on mental simulation when evaluating their own outputs at inference time---and SolidCoder's live execution provides additional value by replacing this unreliable self-assessment with concrete feedback. Overall, these results validate our core hypothesis: the Mental-Reality Gap is a fundamental limitation of mental simulation, not an artifact of specific model weaknesses, and execution-grounded verification provides consistent improvements across model families and training paradigms.

\subsection{Ablation Study}

% ============================================================
% [실험 메모] Ablation Study - Leave-One-Out 방식
% ============================================================
% 목적: 각 S.O.L.I.D. 컴포넌트가 빠졌을 때 성능이 얼마나 떨어지는지 측정
% 벤치마크: CodeContests (GPT-4o) - 개선폭 +4.4%p로 ablation 분석에 적합
% 이유: HumanEval은 +0.6%p (1문제)로 차이가 너무 작아 ablation 의미 없음
%
% ====== 실험 명령어 (main.py 플래그와 일치) ======
% Full:      --enable_shift_left --enable_oracle_assert --enable_live_verify --enable_inter_sim --enable_defensive_test
% w/o S:     --enable_oracle_assert --enable_live_verify --enable_inter_sim --enable_defensive_test
% w/o O:     --enable_shift_left --enable_live_verify --enable_inter_sim --enable_defensive_test
% w/o L:     --enable_shift_left --enable_oracle_assert --enable_inter_sim --enable_defensive_test
% w/o I:     --enable_shift_left --enable_oracle_assert --enable_live_verify --enable_defensive_test
% w/o D:     --enable_shift_left --enable_oracle_assert --enable_live_verify --enable_inter_sim
%
% 결과 (CC, GPT-4o):
% - Full: 77.0% (127/165)
% - w/o S: 53.3% (88/165, -23.7%p) - 가장 큰 하락!
% - w/o I: 64.0% (103/161, -13.0%p)
% - w/o O: 65.4% (108/165, -11.6%p)
% - w/o L: 69.1% (114/165, -7.9%p)
% - w/o D: 70.3% (116/165, -6.7%p)
% - CodeSIM: 72.7%
% ============================================================

% 각 S.O.L.I.D. 컴포넌트의 기여도를 정량화하기 위해, 기존 연구와의 비교를 위한 주요 모델인 GPT-4o를 사용하여 CodeContests 벤치마크에서 leave-one-out ablation study를 수행한다. HumanEval은 현재 LLM 성능이 포화 상태에 가까워(95%+) 미세한 변화를 관찰하기 어려울 것으로 예상되고, APPS는 반대로 난이도가 너무 높아 컴포넌트 간 차이가 낮은 점수에 가려질 수 있다. CodeContests는 적절한 난이도를 보이면서도 각 컴포넌트의 영향이 뚜렷이 드러날 것으로 기대되는 벤치마크이다.
To quantify the contribution of each S.O.L.I.D. component, we conduct a leave-one-out ablation study on CodeContests using GPT-4o, our primary model for comparison with prior work.
We select this benchmark because it offers an appropriate difficulty level where component contributions are likely to be clearly distinguishable.
On HumanEval, current LLMs approach saturation (95\%+), which would make subtle variations difficult to observe reliably.
Conversely, APPS poses such high difficulty that component differences may be obscured by uniformly low scores.
CodeContests strikes a balance, providing sufficient headroom to reveal each component's contribution.

\begin{table}[t]
    \centering
    \small
    \begin{tabular}{@{}lcc@{}}
        \toprule
        \textbf{Configuration} & \textbf{Pass@1} & \textbf{$\Delta$} \\
        \midrule
        GPT-4o Direct  & 42.4 & \textcolor{CrossRed}{$-$34.6} \\
        \addlinespace[3pt]
        \quad w/o \textbf{S}hift-left Planning & 53.3 & \textcolor{CrossRed}{$-$23.7} \\
        \quad w/o \textbf{I}ntermediate Simulation & 64.0 & \textcolor{CrossRed}{$-$13.0} \\
        \quad w/o \textbf{O}racle-based Assertions & 65.4 & \textcolor{CrossRed}{$-$11.6} \\
        \quad w/o \textbf{L}ive Execution & 69.1 & \textcolor{CrossRed}{$-$7.9} \\
        \quad w/o \textbf{D}efensive Accumulation & 70.3 & \textcolor{CrossRed}{$-$6.7} \\
        \midrule
        \textbf{GPT-4o SolidCoder (Full)} & \textbf{77.0} & \textcolor{gray}{--} \\
        \bottomrule
    \end{tabular}
    % {CodeContests(GPT-4o)에 대한 어블레이션. $\Delta$는 전체 버전 대비 감소량을 나타냅니다. 구성 요소는 영향도 순으로 순위가 매겨져 있습니다.}
    \caption{Ablation on CodeContests (GPT-4o). $\Delta$ indicates drop from Full. Components ranked by impact.}
    \label{tab:ablation}
\end{table}


\begin{table*}[t]
    \centering
    \small
    \setlength{\tabcolsep}{3pt}
    \begin{tabular}{@{}lccccccccc@{}}
        \toprule
        & \multicolumn{3}{c}{\textbf{GPT-4o}$^*$} & \multicolumn{3}{c}{\textbf{GPT-OSS-120B}} & \multicolumn{3}{c}{\textbf{Grok-4.1-Fast}} \\
        \cmidrule(lr){2-4} \cmidrule(lr){5-7} \cmidrule(lr){8-10}
        & MapCoder & CodeSIM & SolidCoder & MapCoder & CodeSIM & SolidCoder & MapCoder & CodeSIM & SolidCoder \\
        \midrule
        \rowcolor{green!10}
        \multicolumn{10}{c}{\textbf{HumanEval} \textcolor{green!60!black}{\scriptsize(EASY)}} \\
        Pass@1 (\%) & 90.2 & 95.1 & \textbf{95.7} & 61.0 & \textbf{98.2} & \textbf{98.2} & 95.7 & 97.0 & \textbf{97.6} \\
        API Calls & 9 & \underline{4} & 8 & 6 & \underline{3} & 19 & 8 & \underline{4} & 12 \\
        Tokens (K) & 6.6 & \underline{3.8} & 7.5 & 10.0 & \underline{5.6} & 40.9 & 18.7 & \underline{14.3} & 47.8 \\
        \addlinespace[3pt]
        \rowcolor{yellow!15}
        \multicolumn{10}{c}{\textbf{CodeContests} \textcolor{yellow!60!black}{\scriptsize(MEDIUM)}} \\
        Pass@1 (\%) & 68.3 & 72.7 & \textbf{77.0} & 43.5 & 87.6 & \textbf{91.9} & 85.7 & 95.0 & \textbf{98.1} \\
        API Calls & 11 & \underline{11} & 16 & \underline{4} & 6 & 15 & 8 & \underline{4} & 10 \\
        Tokens (K) & 20.9 & \underline{20.9} & 31.9 & \underline{11.5} & 28.3 & 58.3 & 44.2 & \underline{35.0} & 96.7 \\
        \addlinespace[3pt]
        \rowcolor{red!10}
        \multicolumn{10}{c}{\textbf{APPS} \textcolor{red!60!black}{\scriptsize(HARD)}} \\
        Pass@1 (\%) & 20.7 & 23.3 & \textbf{26.7} & 23.5 & 38.7 & \textbf{40.0} & 33.6 & 41.3 & \textbf{42.0} \\
        API Calls & \underline{19} & 26 & 35 & \underline{9} & 21 & 48 & 15 & \underline{20} & 47 \\
        Tokens (K) & \underline{36.4} & 49.3 & 60.4 & \underline{31.3} & 73.3 & 177.4 & 97.5 & \underline{266.5} & 520.9 \\
        \bottomrule
    \end{tabular}
    % 효율성 비교. 밑줄=가장 효율적, 볼드=최고 Pass@1.
    \caption{Efficiency comparison across models. \underline{Underline}: most efficient; \textbf{bold}: highest Pass@1 per model-benchmark. GPT-OSS-120B shows degraded MapCoder performance due to XML parsing incompatibility. $^*$Primary model for comparison with prior work.}
    \label{tab:efficiency}
\end{table*}


% 표~\ref{tab:ablation}은 각 컴포넌트 제거 시의 성능 변화를 보여준다. 결과는 모든 컴포넌트가 유의미하게 기여함을 확인해준다.
Table~\ref{tab:ablation} presents the results. The findings confirm that all five components contribute significantly to SolidCoder's performance.

% Shift-left Planning (S)가 가장 큰 영향을 미치며, 제거 시 -23.7%p의 급격한 하락을 보인다---다른 어떤 컴포넌트보다 훨씬 큰 하락폭이다. 이는 코드 생성 이전에 edge case를 예측하는 것이 필수적임을 입증한다.
\textbf{Shift-left Planning (S)} emerges as the most critical component, with its removal causing a dramatic $-23.7$\%p drop---by far the largest among all components. This confirms that anticipating edge cases \emph{before} code generation is essential; without it, the model frequently produces solutions that fail on boundary conditions.

% Intermediate Simulation (I)이 두 번째로 큰 기여(-13.0%p)를 보여, mental simulation이 live execution 이전의 보완적 필터로서 논리적 오류를 조기에 포착하는 데 가치를 제공한다.
\textbf{Intermediate Simulation (I)} ranks second ($-13.0$\%p), demonstrating that mental simulation---when strategically positioned as a complementary filter before live execution---provides substantial value in catching logical errors early. 

% Oracle-based Assertions (O)는 -11.6%p의 기여를 보이며, 문제 제약 조건에 기반한 테스트 생성이 정확성 향상에 중요함을 입증한다.
\textbf{Oracle-based Assertions (O)} contributes $-11.6$\%p, showing that grounding test generation in problem constraints improves solution correctness.

% Live Execution (L)은 -7.9%p의 영향을 보여, 순수한 추론만으로는 놓치는 오류를 구체적 실행 피드백이 잡아냄을 검증한다.
\textbf{Live Execution (L)} shows $-7.9$\%p impact, validating that concrete execution feedback catches errors that pure reasoning misses. 

% Defensive Accumulation (D)는 -6.7%p의 기여를 보이며, 실패한 케이스로부터 학습하는 것이 반복적 개선에 가치가 있음을 나타낸다.
\textbf{Defensive Accumulation (D)} contributes $-6.7$\%p, indicating value in learning from failed cases across iterations.

% 종합하면, 사전 계획(S)이 오류를 원천에서 방지하여 가장 큰 효과를 보이고, 시뮬레이션(I)과 실행(L) 기반 검증이 상호 보완적으로 다른 유형의 오류를 포착하며, Oracle(O)과 Defensive(D)가 정제 과정 전반에 걸쳐 견고한 테스트 커버리지를 보장한다.
Overall, these results demonstrate that the five S.O.L.I.D. components work in concert: proactive edge-case planning (S) yields the largest gains by preventing errors at the source, simulation-based (I) and execution-based (L) verification catch complementary error types, and oracle-grounded assertions (O) with defensive accumulation (D) ensure robust test coverage throughout the refinement process.

\subsection{Efficiency Analysis}

% ============================================================
% [효율성 분석] - 난이도 기반 선택적 적용 관점
% ============================================================
% 핵심 메시지:
% - SolidCoder는 토큰을 더 소비하지만, 어려운 문제에서 효율적인 trade-off
% - 문제가 어려워질수록 토큰 대비 정확도 개선 효율이 높아짐
% - 향후 난이도 판단 → 방법론 선택적 적용 가능
% ============================================================

% SolidCoder의 실행 기반 검증은 추가 비용을 발생시킵니다. 표~\ref{tab:efficiency}에서 모델별 효율성을 비교합니다.
SolidCoder's execution-grounded verification introduces additional cost. Table~\ref{tab:efficiency} compares efficiency across models.

% HHumanEval과 같은 쉬운 문제에서는 SolidCoder의 오버헤드가 비효율적입니다. GPT-4o는 +0.6%p의 성능 향상을 위해 +97% 더 많은 토큰을 사용합니다. 그러나 기본 LLM 모델이 어려움을 겪는 CodeContests 및 APPS와 같은 문제에서는 추가 비용이 비례적으로 성능 향상을 가져옵니다. CodeContests에서는 +4.3%p의 성능 향상을 위해 +50%의 토큰을 사용했고, APPS에서는 +3.4%p의 성능 향상을 위해 +23%의 토큰을 사용했습니다. 강화 학습 후 모델도 비슷한 패턴을 보이며, 더 높은 토큰 오버헤드에도 불구하고 CodeContests에서 일관된 성능 향상(GPT-OSS +4.3%p, Grok +3.1%p)을 달성합니다.
On easy problems like HumanEval, SolidCoder's overhead is inefficient: GPT-4o uses +97\% more tokens for only +0.6\%p gain. However, on problems where base LLMs struggle---CodeContests and APPS---the additional cost yields proportional improvements: +50\% tokens for +4.3\%p on CodeContests, +23\% for +3.4\%p on APPS. RL post-trained models show similar patterns, achieving consistent gains on CodeContests (+4.3\%p for GPT-OSS, +3.1\%p for Grok) despite higher token overhead.

% 이는 난이도를 고려한 라우팅이 SolidCoder의 가치를 극대화할 수 있음을 시사합니다. 즉, 간단한 문제에는 경량화된 방법을 적용하고, 머릿속으로 시뮬레이션하는 것이 실패할 가능성이 가장 높은 복잡한 작업에는 실행 기반의 완벽한 검증을 적용하는 것입니다.
This suggests that difficulty-aware routing could maximize SolidCoder's value: applying lightweight methods to simple problems while reserving full execution-grounded verification for complex tasks where mental simulation is most likely to fail.

\subsection{Qualitative Example}
% SolidCoder가 Mental-Reality Gap을 어떻게 해결하는지 보여주기 위해 정성 분석을 수행합니다.
% 그림 2는 핵심 차이점을 보여주는 개념적 예시입니다: mental simulation은 예상 출력을 환각하여 결함 있는 코드를 잘못 검증할 수 있지만, SolidCoder의 Live Execution은 구체적인 테스트 실패로 이러한 버그를 포착합니다.
% 부록에서는 실제 실험에서 추출한 전체 실행 로그를 제공합니다.
Figure~\ref{fig:2} presents an illustrative example of how SolidCoder addresses the Mental-Reality Gap: mental simulation may incorrectly validate flawed code by hallucinating expected outputs, while SolidCoder's Live Execution catches such bugs through concrete test failure. A complete execution trace from our experiments, including full LLM inputs and outputs, is provided in Appendix~\ref{app:example}.

\section{Conclusion and Future Work}
% LLM이 실행 추적을 환각하며 버그가 있는 코드가 제대로 작동한다고 스스로 확신하는 정신-현실 간극(Mental-Reality Gap)에 착안하여, 구체적인 실행을 통해 이를 해소하기 위해 S.O.L.I.D. 아키텍처를 기반으로 하는 SolidCoder를 제안했습니다.
% 핵심적인 통찰은 간단합니다: 상상하지 말고 실행하세요.
% GPT-4o 실험 결과, SolidCoder는 HumanEval 95.7%(+0.6%p), CodeContests 77.0%(+4.3%p), APPS 26.7%(+3.4%p)를 달성했습니다.
% 이러한 개선은 RL post-trained 모델에서도 일관되게 나타나며, 특히 CodeContests에서 가장 큰 개선을 보입니다.
% 이는 런타임 현실에 기반한 검증이 견고한 코드 합성에 필수적임을 보여줍니다.
% 향후 연구에서는 난이도 기반 라우팅과 Mental-Reality Gap의 근본 원인을 탐구할 것입니다.
Motivated by the \textbf{Mental-Reality Gap}---where LLMs hallucinate execution traces, convincing themselves that buggy code works correctly---we proposed \textbf{SolidCoder} with the \textbf{S.O.L.I.D.} architecture to bridge it through concrete execution. Our core insight is simple: \textit{don't imagine---execute}. Experiments with GPT-4o demonstrate new state-of-the-art pass@1 results: \textbf{95.7\%} on HumanEval, \textbf{77.0\%} on CodeContests, and \textbf{26.7\%} on APPS---improving over CodeSIM by +0.6\%p, +4.3\%p, and +3.4\%p respectively. These improvements generalize consistently across RL post-trained models, with the largest gains on CodeContests where mental simulation becomes unreliable. This validates that grounding verification in runtime reality is essential for robust code synthesis. Future work will explore difficulty-aware routing to maximize efficiency, and investigate the fundamental causes of when and why the Mental-Reality Gap manifests.

\section*{Limitations}

% 저희 연구는 향후 연구 방향을 시사하는 몇 가지 한계점을 가지고 있습니다.
Our work has several limitations that point to future research directions.

% 첫째, Live Execution은 현재 Python만 지원합니다. S.O.L.I.D. 아키텍처 자체는 언어에 구애받지 않지만, 다른 언어로 확장하려면 언어별 샌드박싱을 구현해야 합니다. 이는 근본적인 한계라기보다는 다루기 쉬운 엔지니어링 작업입니다.
\textbf{First}, Live Execution currently supports \textbf{Python only}. The S.O.L.I.D. architecture itself is language-agnostic, but extending to other languages requires implementing language-specific sandboxing---a tractable engineering effort rather than a fundamental limitation.

% 둘째, Oracle 기반 어설션과 Shift-left Planning은 심적 시뮬레이션에 대한 의존도를 낮추지만, 여전히 LLM의 유효한 속성과 예외 사례를 식별하는 능력에 의존합니다. 속성 검증(예: ``출력 길이가 입력 길이와 같음'')은 정확한 출력 예측보다 본질적으로 제약이 많아 환각 위험을 줄이지만, 환각 속성은 여전히 존재할 수 있다고 주장합니다. 이는 모든 LLM 기반 테스트 생성 방식이 공유하는 과제입니다.
\textbf{Second}, while Oracle-based Assertions and Shift-left Planning reduce reliance on mental simulation, they still depend on the LLM's ability to identify valid properties and edge cases. We argue that property verification (e.g., ``output length equals input length'') is inherently more constrained than exact output prediction, reducing hallucination risk, but \textbf{hallucinated properties remain possible}---a challenge shared by all LLM-based test generation approaches.

% 셋째, 본 평가는 정신-현실 간 격차를 명확하게 분리하는 기능 수준 벤치마크에 중점을 둡니다. 다중 파일 조정(예: SWE-bench)을 포함하는 저장소 수준 작업에 대한 일반화는 아직 검증되지 않았습니다. 그러나 S.O.L.I.D. 구성 요소, 특히 실시간 실행 및 방어적 누적은 기존 저장소 수준 프레임워크를 보완할 수 있다고 생각합니다.
\textbf{Third}, our evaluation focuses on \textbf{function-level benchmarks} that cleanly isolate the Mental-Reality Gap. Generalization to repository-level tasks involving multi-file coordination (e.g., SWE-bench) remains to be validated; however, we believe S.O.L.I.D. components---particularly Live Execution and Defensive Accumulation---could complement existing repository-level frameworks.

% 넷째, SolidCoder는 정신-현실 간 격차를 줄이지만 완전히 없애지는 않습니다. 중급 시뮬레이션은 여전히 정신 추론을 사용하며, Shift-left 커버리지는 LLM 이해도에 따라 제한됩니다. 그럼에도 불구하고, 본 실험은 구체적인 실행에 대한 부분적인 기초 지식만으로도 순수한 정신 시뮬레이션보다 상당한 정확성 향상을 가져온다는 것을 보여줍니다.
\textbf{Fourth}, SolidCoder \textbf{reduces but does not eliminate} the Mental-Reality Gap. Intermediate Simulation still employs mental reasoning, and Shift-left coverage is bounded by LLM understanding. Nevertheless, our experiments demonstrate that even partial grounding in concrete execution yields substantial correctness gains over pure mental simulation.

% 다섯째, 모델 범위 제한. 기존 연구(CodeSIM)가 GPT-3.5, GPT-4, GPT-4o를 사용한 것에서 나아가, 본 연구는 GPT-4o, GPT-OSS-120B(오픈소스 GPT), Grok-4.1-Fast(비-GPT 프론티어)로 확장하여 RL post-trained 모델까지 평가 범위를 넓혔다. 그러나 실험 비용과 시간의 제약으로 Gemini, Claude 등 다른 강력한 모델 패밀리는 포함하지 못했다. 향후 연구에서 이러한 모델들로 일반화 가능성을 검증할 수 있을 것이다.
\textbf{Fifth}, our \textbf{model coverage}, while broader than prior work, remains incomplete. CodeSIM evaluated on GPT-3.5, GPT-4, and GPT-4o; we extend this to include GPT-OSS-120B (an open-source GPT variant) and Grok-4.1-Fast (a non-GPT frontier model), thereby covering RL post-trained models. However, due to experimental cost and time constraints, we did not evaluate other strong model families such as Gemini and Claude. Future work could validate generalization across these additional architectures.

% 여섯째, 벤치마크 오염 가능성. HumanEval(2021), CodeContests(2022), APPS(2021)는 현재 평가에 사용된 모델들의 학습 데이터에 포함되었을 가능성을 배제할 수 없다. 그러나 본 연구의 핵심은 절대 성능이 아닌 동일 조건에서의 방법론 간 비교이다. 모든 전략(Direct, CoT, MapCoder, CodeSIM, SolidCoder)이 동일한 모델과 벤치마크에서 평가되었으므로, 관찰된 상대적 개선은 벤치마크 오염과 무관하게 유효하다.
\textbf{Sixth}, \textbf{benchmark contamination} cannot be ruled out. HumanEval \citep{chen2021evaluating}, CodeContests \citep{li2022competition}, and APPS \citep{hendrycks2021measuring} may have been included in the training data of the models we evaluate. However, the core contribution of this work lies in \textit{relative comparisons} under identical conditions, not absolute performance claims. Since all strategies---Direct, CoT, MapCoder, CodeSIM, and SolidCoder---are evaluated on the same models and benchmarks, the observed relative improvements remain valid regardless of potential contamination.

\bibliography{custom}

\appendix

% ============================================================
% APPENDIX 구조
% ============================================================
% SolidCoder Appendix 구성:
%   A. Algorithm Details (SolidCoder 워크플로 알고리즘)
%   B. S.O.L.I.D. Component Prompts (각 컴포넌트별 프롬프트)
%   C. Live Execution Safety (샌드박스 구현 세부사항)
%   D. Complete Example (전체 실행 trace 예시)
% ============================================================

\section{Algorithm Details}
\label{app:algorithm}

% SolidCoder의 전체 워크플로를 Algorithm~\ref{alg:solidcoder}에 제시합니다. 알고리즘은 세 가지 주요 단계로 구성됩니다: (1) Shift-left Planning을 통한 edge case 식별 및 계획 생성, (2) Assumption Breaking Phase에서 Oracle 테스트와 Live Execution을 통한 검증, (3) Traditional Debugging Phase에서의 반복적 코드 수정. 각 S.O.L.I.D. 컴포넌트는 파란색으로 표시되어 있습니다.
Algorithm~\ref{alg:solidcoder} presents the complete SolidCoder workflow. The algorithm consists of three main stages: (1) edge case identification and plan generation via Shift-left Planning, (2) verification through Oracle tests and Live Execution in the Assumption Breaking Phase, and (3) iterative code refinement in the Traditional Debugging Phase. Each S.O.L.I.D. component is marked in blue.

\begin{algorithm}[h]
\caption{The SolidCoder Workflow}
\label{alg:solidcoder}
\small
\begin{algorithmic}[1]
\Require Problem $P$, Sample I/O $T$, max\_plan $p$, max\_debug $d$, max\_assume $a$
\Ensure Solution Code $S$
\State $accumulated \leftarrow \emptyset$
\For{$i = 1$ to $p$}
    \State \textcolor{blue}{\textbf{[S]}} $edges \leftarrow \Call{IdentifyEdgeCases}{P}$
    \State $plan \leftarrow \Call{GeneratePlan}{P, edges}$
    \State $valid \leftarrow \Call{SimulatePlan}{P, plan}$ \Comment{Mental Verification}
    \If{$\neg valid$}
        \State $plan \leftarrow \Call{RefinePlan}{P, plan}$
    \EndIf
    \State $S \leftarrow \Call{GenerateCode}{P, plan}$
    \State \textcolor{blue}{\textbf{[I]}} $S \leftarrow \Call{IntermediateSimulation}{P, S}$ \Comment{Complementary filter}

    \For{$k = 1$ to $a$} \Comment{Assumption Breaking Phase}
        \State \textcolor{blue}{\textbf{[O]}} $test \leftarrow \Call{GenerateOracleTest}{P, S}$
        \State \textcolor{blue}{\textbf{[L]}} $result \leftarrow \Call{LiveExecute}{S, test}$
        \If{$result = \texttt{FAIL}$}
            \State \textcolor{blue}{\textbf{[D]}} $accumulated \leftarrow accumulated \cup \{test\}$
            \State $S \leftarrow \Call{FixCode}{P, S, test}$
        \Else
            \State \textbf{break} \Comment{Verification Passed}
        \EndIf
    \EndFor

    \If{$\Call{PassAll}{S, T \cup accumulated}$}
        \State \Return $S$
    \EndIf

    \For{$j = 1$ to $d$} \Comment{Traditional Debugging Phase}
        \State $S \leftarrow \Call{DebugCode}{P, plan, S}$
        \If{$\Call{PassAll}{S, T \cup accumulated}$}
            \State \Return $S$
        \EndIf
    \EndFor
\EndFor
\State \Return $S$
\end{algorithmic}
\end{algorithm}

\section{S.O.L.I.D. Component Prompts}
\label{app:prompts}

This appendix presents the complete prompts used in SolidCoder's S.O.L.I.D. components. Variables in curly braces (e.g., \texttt{\{problem\}}) are replaced with actual values at runtime. Each component is color-coded to match Figure~\ref{fig:1}.

\subsection{[S] Shift-left Planning Prompt}

The Shift-left prompt extends standard planning by requiring edge case identification \textit{before} algorithm design:

\begin{promptbox}[title={\textcolor{white}{[S] Shift-left Planning Prompt}}]{ShiftColor}
\small\ttfamily
You are an expert software engineer.\\
You are given a problem and you need to generate a detailed plan to solve it.\\[0.5em]
\#\# Problem\\
\{problem\}\\[0.5em]
\#\# Killer Edge Cases\\
Identify 3 potential edge cases that could break a naive solution:\\
1. Empty/minimal input\\
2. Maximum constraint input\\
3. Special pattern (all same, alternating, etc.) or Boundary values\\[0.5em]
\#\# Plan\\
Write a detailed plan that handles these edge cases.\\
The plan should be step-by-step and easy to implement.
\end{promptbox}

\subsection{[O] Oracle-based Assertion Prompt}

The Oracle prompt instructs the LLM to act as a ``Red Team'' tester, generating property-based assertions rather than exact output predictions:

\begin{promptbox}[title={\textcolor{white}{[O] Oracle-based Assertion Prompt}}]{OracleColor}
\small\ttfamily
You are an expert software tester (Red Team).\\
Your goal is to break the following code by finding hidden assumptions.\\[0.5em]
\#\# Problem\\
\{problem\}\\[0.5em]
\#\# Code\\
```\{language\}\\
\{code\}\\
```\\[0.5em]
\#\# Your Task\\
1. Identify a weak assumption (Type, Value, Structure, or Relationship).\\
2. Produce a **Python test script** that calls the target function with a breaking input.\\
3. Include an assert that would fail if the assumption is violated.\\[0.5em]
Format:\\
Assumption: <short text>\\
Test Script:\\
```python\\
\# call the function defined above\\
result = <call>\\
assert <oracle about result>\\
```
\end{promptbox}

\subsection{[L] Live Execution}

Live Execution does not use a separate prompt---it directly executes the test scripts generated by the Oracle component in a sandboxed Python environment. The execution wrapper is shown below:

\begin{codebox}[title={\textcolor{white}{[L] Live Execution Wrapper}}, colframe=LiveColor!80!black, colback=LiveColor!5!white, boxed title style={colback=LiveColor!80!black}]
\small\ttfamily
def \_concrete\_verify\_script(self, code, test\_script, timeout=5):\\
\hspace*{1em}full\_code = f\textquotedbl\textquotedbl\textquotedbl\\
\hspace*{1em}import sys\\
\hspace*{1em}import math\\
\hspace*{1em}from typing import List, Dict, Any, Optional, Union, Tuple\\[0.3em]
\hspace*{1em}\{code\}\\[0.3em]
\hspace*{1em}\{test\_script\}\\
\hspace*{1em}\textquotedbl\textquotedbl\textquotedbl\\
\hspace*{1em}try:\\
\hspace*{2em}safe\_builtins = builtins.\_\_dict\_\_.copy()\\
\hspace*{2em}safe\_builtins["input"] = lambda *\_, **\_\_: $\backslash$\\
\hspace*{3em}(\_ for \_ in ()).throw(RuntimeError("input() disabled"))\\
\hspace*{2em}exec\_globals = \{"\_\_builtins\_\_": safe\_builtins\}\\
\hspace*{2em}function\_with\_timeout(exec, (full\_code, exec\_globals), timeout)\\
\hspace*{2em}return "PASS"\\
\hspace*{1em}except AssertionError:\\
\hspace*{2em}return "FAIL\_ASSERT"\\
\hspace*{1em}except Exception:\\
\hspace*{2em}return "FAIL\_CRASH"
\end{codebox}

\noindent See Appendix~\ref{app:safety} for detailed safety mechanisms.

\subsection{[I] Intermediate Simulation Prompt}

The Intermediate Simulation prompt asks the LLM to mentally trace code execution before live verification:

\begin{promptbox}[title={\textcolor{white}{[I] Intermediate Simulation Prompt}}]{InterColor}
\small\ttfamily
\#\# Problem\\
\{problem\}\\[0.5em]
\#\# Code\\
```\{language\}\\
\{code\}\\
```\\[0.5em]
\#\# Your Task: Mental Execution\\
1. Select a sample input.\\
2. Trace the code execution step-by-step.\\
3. Track variable values.\\
4. Predict the final output.\\[0.5em]
If you find a logic error or mismatch with the plan, output: **CODE\_SIMULATION\_FAILED**\\
If the code seems correct, output: **CODE\_SIMULATION\_PASSED**
\end{promptbox}

\subsection{[D] Defensive Accumulation}

Defensive Accumulation is implemented as a runtime mechanism rather than a prompt. When Live Execution detects a failing test case, the test is automatically added to an \texttt{accumulated\_tests} collection. All subsequent code modifications must pass \textit{every} accumulated test before being accepted:

\begin{codebox}[title={\textcolor{white}{[D] Defensive Accumulation Logic}}, colframe=DefensiveColor!80!black, colback=DefensiveColor!5!white, boxed title style={colback=DefensiveColor!80!black}]
\small\ttfamily
\textcolor{gray}{\# In the assumption breaking loop:}\\
verify\_status = self.\_concrete\_verify\_script(code, test\_script)\\[0.3em]
if verify\_status in \{"FAIL\_CRASH", "FAIL\_ASSERT"\}:\\
\hspace*{1em}\textcolor{gray}{\# [D] Defensive Accumulation: store failing test}\\
\hspace*{1em}if self.enable\_defensive\_test:\\
\hspace*{2em}self.accumulated\_inputs.append(test\_script)\\
\hspace*{2em}additional\_io.append(test\_script)\\[0.3em]
\hspace*{1em}\textcolor{gray}{\# Fix the code based on the failing test}\\
\hspace*{1em}code = self.\_fix\_with\_test\_script(problem, code, test\_script)\\
\hspace*{1em}continue \textcolor{gray}{\# Re-verify with accumulated tests}
\end{codebox}

\subsection{Test Case Validation (Judge Prompt)}

To prevent false positives from invalid test cases, SolidCoder includes a Judge component that validates generated tests:

\begin{promptbox}[title={\textcolor{white}{Judge Prompt}}, colframe=gray!70!black, colback=gray!5!white, boxed title style={colback=gray!70!black}]{gray}
\small\ttfamily
You are an impartial Judge.\\[0.5em]
\#\# Problem\\
\{problem\}\\[0.5em]
\#\# Proposed Test Case\\
```python\\
\{test\_script\}\\
```\\[0.5em]
\#\# Task\\
Determine if this test case is VALID and CORRECT for the given problem.\\
1. Does the input satisfy all constraints (e.g. range, type, format)?\\
2. Is the asserted output (if any) logically correct?\\
3. Is it a fair test?\\[0.5em]
If the test case is valid, output: **VALID**\\
If it violates constraints or expects wrong output, output: **INVALID**
\end{promptbox}

\subsection{Code Fix Prompt}

When a vulnerability is found, SolidCoder uses this prompt to fix the code:

\begin{promptbox}[title={\textcolor{white}{Code Fix Prompt}}, colframe=gray!70!black, colback=gray!5!white, boxed title style={colback=gray!70!black}]{gray}
\small\ttfamily
\#\# Problem\\
\{problem\}\\[0.5em]
\#\# Code\\
```\{language\}\\
\{code\}\\
```\\[0.5em]
\#\# Vulnerability Found\\
The following test script failed (Crash or Assertion Error):\\
```python\\
\{test\_script\}\\
```\\[0.5em]
\#\# Task\\
Fix the code to handle this edge case.\\
Wrap the code in ```\{language\} ... ``` block.
\end{promptbox}

\section{Live Execution Safety Mechanisms}
\label{app:safety}

SolidCoder implements several safety mechanisms for Live Execution to prevent malicious or buggy code from affecting the host system. Table~\ref{tab:safety} summarizes these mechanisms. Full implementation details are provided in the \textcolor[rgb]{0.6, 0.6, 0.0}{[L] Live Execution} code listing in Appendix~\ref{app:prompts}.

\begin{table*}[t]
\centering
\small
\renewcommand{\arraystretch}{1.3}
\begin{tabular}{@{}>{\raggedright\arraybackslash}p{2.8cm}>{\raggedright\arraybackslash}p{4.2cm}>{\raggedright\arraybackslash}p{8.5cm}@{}}
\toprule
\rowcolor{SafetyColor!15}
\textbf{Mechanism} & \textbf{Implementation} & \textbf{Purpose \& Rationale} \\
\midrule
Timeout & 5-second limit per test via \texttt{function\_with\_timeout} & Prevents infinite loops from blocking the pipeline. Correct solutions complete in milliseconds; 5s provides margin while catching runaway executions. \\
\arrayrulecolor{gray!30}\hline\arrayrulecolor{black}
Input Blocking & \texttt{input()} $\rightarrow$ \texttt{RuntimeError} stub & Prevents scripts from hanging indefinitely waiting for user input that will never arrive in automated execution. \\
\arrayrulecolor{gray!30}\hline\arrayrulecolor{black}
Isolated Namespace & Fresh \texttt{exec\_globals} per execution & Prevents test scripts from interfering with each other or accumulating state. Each test starts with a clean environment. \\
\arrayrulecolor{gray!30}\hline\arrayrulecolor{black}
Exception Classification & \texttt{PASS}, \texttt{FAIL\_ASSERT}, \texttt{FAIL\_CRASH} & Distinguishes logical failures (\texttt{AssertionError}) from runtime bugs. Both trigger Defensive Accumulation with different diagnostic signals. \\
\bottomrule
\end{tabular}
\caption{Safety mechanisms implemented in SolidCoder's Live Execution component to ensure reliable and secure code evaluation.}
\label{tab:safety}
\end{table*}

% Complete Example 섹션은 전체 페이지 너비를 사용합니다.
\onecolumn

\section{Complete Example}
\label{app:example}

% 이 섹션은 SolidCoder가 HumanEval 문제를 푸는 전체 과정을 보여줍니다.
% GPT-4o를 사용한 실제 실험에서 추출한 실행 로그입니다.

This section presents a complete execution trace of SolidCoder solving HumanEval problem \#0 (\texttt{has\_close\_elements}). The trace demonstrates the full S.O.L.I.D. pipeline: \textbf{[S]} Shift-left Planning identifies edge cases before generating the plan; \textbf{[I]} Intermediate Simulation traces the generated code mentally; \textbf{[O]} Oracle-based Assumption Breaking generates adversarial test cases; and \textbf{[L]} Live Execution validates correctness through actual runtime. In this example, the code passes on the first attempt, so \textbf{[D]} Defensive Accumulation is not triggered. All prompts and LLM responses are shown verbatim from the experiment log.

\vspace{3mm}

\begin{tcolorbox}[breakable, colback=white, colframe=black!70, boxrule=0.5pt, arc=2pt, left=4mm, right=4mm, top=3mm, bottom=3mm]

\vspace{2mm}
\textbf{\large Complete Execution Trace: HumanEval/0 (\texttt{has\_close\_elements}) with GPT-4o}
\hrule
\vspace{5mm}

%==============================================================================
\noindent\textbf{\Large Problem}
%==============================================================================

\begin{verbatim}
from typing import List

def has_close_elements(numbers: List[float], threshold: float) -> bool:
    """ Check if in given list of numbers, are any two numbers closer to each 
    other than given threshold.
    >>> has_close_elements([1.0, 2.0, 3.0], 0.5)
    False
    >>> has_close_elements([1.0, 2.8, 3.0, 4.0, 5.0, 2.0], 0.3)
    True
    """
\end{verbatim}

\vspace{0.5cm} \hrule \vspace{0.5cm}

%==============================================================================
\noindent\textbf{\Large [S] Shift-left Planning}
%==============================================================================

\textbf{Prompt:}
\begin{verbatim}
You are an expert software engineer.
You are given a problem and you need to generate a detailed plan to solve it.

## Problem
[problem code shown above]

## Killer Edge Cases
Identify 3 potential edge cases that could break a naive solution:
1. Empty/minimal input
2. Maximum constraint input
3. Special pattern (all same, alternating, etc.) or Boundary values

## Plan
Write a detailed plan that handles these edge cases.
The plan should be step-by-step and easy to implement.
\end{verbatim}

\textbf{LLM Response:}

To solve the problem of determining if any two numbers in a list are closer to each other than a given threshold, we need to consider potential edge cases and develop a robust plan. Here's a detailed plan to address the problem and handle edge cases:

\textbf{\#\#\# Edge Cases}

\textbf{1. Empty/Minimal Input:}
\begin{itemize}
    \item An empty list or a list with a single element should return \texttt{False} because there are no pairs of numbers to compare.
\end{itemize}

\textbf{2. Maximum Constraint Input:}
\begin{itemize}
    \item Consider the scenario where the list contains the maximum number of elements allowed by the system's memory. The solution should be efficient in terms of time and space complexity to handle large inputs.
\end{itemize}

\textbf{3. Special Pattern or Boundary Values:}
\begin{itemize}
    \item Lists where all elements are the same, which should return \texttt{True} if the threshold is greater than zero.
    \item Lists with alternating values or values that are very close to each other but not within the threshold.
    \item Lists with very large or very small floating-point numbers to test precision and handling of floating-point arithmetic.
\end{itemize}

\textbf{\#\#\# Plan}

\textbf{1. Initial Checks:}
\begin{itemize}
    \item If the list \texttt{numbers} is empty or contains only one element, return \texttt{False} immediately since no pairs exist.
\end{itemize}

\textbf{2. Sorting:}
\begin{itemize}
    \item Sort the list of numbers. Sorting helps in efficiently finding close elements because if two numbers are close, they will be adjacent in a sorted list.
\end{itemize}

\textbf{3. Iterate and Compare:}
\begin{itemize}
    \item Iterate through the sorted list and compare each pair of adjacent elements.
    \item For each pair \texttt{(numbers[i], numbers[i+1])}, calculate the absolute difference.
    \item If the difference is less than the threshold, return \texttt{True}.
\end{itemize}

\textbf{4. Return Result:}
\begin{itemize}
    \item If no such pair is found after checking all adjacent pairs, return \texttt{False}.
\end{itemize}

\textbf{5. Complexity Consideration:}
\begin{itemize}
    \item Sorting the list takes $O(n \log n)$ time, and iterating through the list takes $O(n)$ time. Thus, the overall time complexity is $O(n \log n)$, which is efficient for large inputs.
\end{itemize}

\textbf{6. Precision Handling:}
\begin{itemize}
    \item Ensure that floating-point arithmetic is handled correctly by using Python's built-in functions which are designed to manage floating-point precision.
\end{itemize}

\vspace{0.5cm} \hrule \vspace{0.5cm}

%==============================================================================
\noindent\textbf{\Large Plan Simulation}
%==============================================================================

\textbf{Prompt:}
\begin{verbatim}
You are a programmer tasked with verifying a plan to solve a given problem 
using the **Python3** programming language.

[Problem and Plan shown above]

**Expected Output:**
Your response must be structured as follows:

### Simulation
- Take a sample input and apply plan step by step to get the output.
- Compare the generated output with the sample output to verify if your plan 
  works as expected.

### Plan Evaluation
- If the simulation is successful write **No Need to Modify Plan**.
- Otherwise write **Plan Modification Needed**.
\end{verbatim}

\textbf{LLM Response:}

\textbf{\#\#\# Simulation}

Let's simulate the plan with the provided test cases to verify its correctness.

\textbf{Test Case 1:} \texttt{numbers = [1.0, 2.0, 3.0]}, \texttt{threshold = 0.5}

\begin{enumerate}
    \item \textbf{Initial Checks:} The list has more than one element, so we proceed.
    \item \textbf{Sorting:} The list is already sorted: \texttt{[1.0, 2.0, 3.0]}.
    \item \textbf{Iterate and Compare:}
    \begin{itemize}
        \item Compare \texttt{1.0} and \texttt{2.0}: $|1.0 - 2.0| = 1.0$, which is not less than \texttt{0.5}.
        \item Compare \texttt{2.0} and \texttt{3.0}: $|2.0 - 3.0| = 1.0$, which is not less than \texttt{0.5}.
    \end{itemize}
    \item \textbf{Return Result:} No pairs found with a difference less than \texttt{0.5}, so return \texttt{False}.
\end{enumerate}

Output: \texttt{False} (Matches expected output)

\textbf{Test Case 2:} \texttt{numbers = [1.0, 2.8, 3.0, 4.0, 5.0, 2.0]}, \texttt{threshold = 0.3}

\begin{enumerate}
    \item \textbf{Initial Checks:} The list has more than one element, so we proceed.
    \item \textbf{Sorting:} Sort the list: \texttt{[1.0, 2.0, 2.8, 3.0, 4.0, 5.0]}.
    \item \textbf{Iterate and Compare:}
    \begin{itemize}
        \item Compare \texttt{1.0} and \texttt{2.0}: $|1.0 - 2.0| = 1.0$, which is not less than \texttt{0.3}.
        \item Compare \texttt{2.0} and \texttt{2.8}: $|2.0 - 2.8| = 0.8$, which is not less than \texttt{0.3}.
        \item Compare \texttt{2.8} and \texttt{3.0}: $|2.8 - 3.0| = 0.2$, which is less than \texttt{0.3}.
    \end{itemize}
    \item \textbf{Return Result:} A pair (\texttt{2.8}, \texttt{3.0}) is found with a difference less than \texttt{0.3}, so return \texttt{True}.
\end{enumerate}

Output: \texttt{True} (Matches expected output)

\textbf{\#\#\# Plan Evaluation}

The simulation of the plan with the provided test cases shows that the implementation works as expected. The plan correctly identifies whether there are any two numbers in the list that are closer to each other than the given threshold.

\textbf{No Need to Modify Plan}

\vspace{0.5cm} \hrule \vspace{0.5cm}

%==============================================================================
\noindent\textbf{\Large Code Generation}
%==============================================================================

\textbf{Prompt:}
\begin{verbatim}
You are a programmer tasked with solving a given problem using the **Python3** 
programming language. See the plan to solve the plan and implement code.

[Problem and Plan shown above]

--------
**Important Instructions:**
- Do not add any explanation.
- The generated **Python3** code must be inside a triple backtick (```) block.
\end{verbatim}

\textbf{LLM Response:}

\begin{verbatim}
from typing import List

def has_close_elements(numbers: List[float], threshold: float) -> bool:
    # Step 1: Initial check for minimal input
    if len(numbers) < 2:
        return False
    
    # Step 2: Sort the list
    numbers.sort()
    
    # Step 3: Iterate and compare adjacent elements
    for i in range(len(numbers) - 1):
        if abs(numbers[i] - numbers[i + 1]) < threshold:
            return True
    
    # Step 4: Return result if no close elements found
    return False
\end{verbatim}

\vspace{0.5cm} \hrule \vspace{0.5cm}

%==============================================================================
\noindent\textbf{\Large [I] Intermediate Simulation}
%==============================================================================

\textbf{Prompt:}
\begin{verbatim}
## Problem
[problem code]

## Code
[generated code]

## Your Task: Mental Execution
1. Select a sample input.
2. Trace the code execution step-by-step.
3. Track variable values.
4. Predict the final output.

If you find a logic error or mismatch with the plan, output: **CODE_SIMULATION_FAILED**
If the code seems correct, output: **CODE_SIMULATION_PASSED**
\end{verbatim}

\textbf{LLM Response:}

Let's select a sample input and trace the code execution step-by-step.

\textbf{Sample Input:}
\begin{verbatim}
numbers = [1.0, 2.8, 3.0, 4.0, 5.0, 2.0]
threshold = 0.3
\end{verbatim}

\textbf{Code Execution:}

\textbf{1. Initial Check:}
\begin{itemize}
    \item The code checks if the length of \texttt{numbers} is less than 2.
    \item \texttt{len(numbers) = 6}, so the check fails, and we proceed to the next step.
\end{itemize}

\textbf{2. Sorting the List:}
\begin{itemize}
    \item The list \texttt{numbers} is sorted.
    \item Sorted \texttt{numbers = [1.0, 2.0, 2.8, 3.0, 4.0, 5.0]}.
\end{itemize}

\textbf{3. Iterate and Compare Adjacent Elements:}

We iterate over the sorted list and compare adjacent elements.

\begin{itemize}
    \item \textbf{Iteration 1} (\texttt{i = 0}):
    \begin{itemize}
        \item Compare \texttt{numbers[0]} and \texttt{numbers[1]}: $|1.0 - 2.0| = 1.0$
        \item $1.0 \geq 0.3$, so continue to the next iteration.
    \end{itemize}
    
    \item \textbf{Iteration 2} (\texttt{i = 1}):
    \begin{itemize}
        \item Compare \texttt{numbers[1]} and \texttt{numbers[2]}: $|2.0 - 2.8| = 0.8$
        \item $0.8 \geq 0.3$, so continue to the next iteration.
    \end{itemize}
    
    \item \textbf{Iteration 3} (\texttt{i = 2}):
    \begin{itemize}
        \item Compare \texttt{numbers[2]} and \texttt{numbers[3]}: $|2.8 - 3.0| = 0.2$
        \item $0.2 < 0.3$, so the condition is satisfied, and the function returns \texttt{True}.
    \end{itemize}
\end{itemize}

\textbf{4. Return Result:}
\begin{itemize}
    \item Since we found two numbers (\texttt{2.8} and \texttt{3.0}) that are closer than the threshold (\texttt{0.3}), the function returns \texttt{True}.
\end{itemize}

\textbf{Conclusion:} The code correctly identifies that there are two numbers in the list that are closer to each other than the given threshold. The logic and execution match the expected behavior.

\textbf{Output: CODE\_SIMULATION\_PASSED}

\vspace{0.5cm} \hrule \vspace{0.5cm}

%==============================================================================
\noindent\textbf{\Large [O] Oracle-based Assumption Breaking (Round 1/3)}
%==============================================================================

\textbf{Prompt:}
\begin{verbatim}
You are an expert software tester (Red Team).
Your goal is to break the following code by finding hidden assumptions.

## Problem
[problem code]

## Code
[generated code]

## Your Task
1. Identify a weak assumption (Type, Value, Structure, or Relationship).
2. Produce a **Python test script** that calls the target function with a 
   breaking input.
3. Include an assert that would fail if the assumption is violated.

Format:
Assumption: <short text>
Test Script:
```python
result = <call>
assert <oracle about result>
```
\end{verbatim}

\textbf{LLM Response:}

\textbf{Assumption:} The function assumes that the input list \texttt{numbers} contains only valid floating-point numbers and does not handle special floating-point values like \texttt{NaN} (Not a Number) or \texttt{inf} (infinity), which can lead to unexpected behavior.

\textbf{Test Script:}
\begin{verbatim}
# call the function defined above
result = has_close_elements([1.0, float('nan'), 3.0], 0.5)
assert result == False, "The function should return False when the list 
contains NaN, as NaN comparisons are not valid."

# call the function defined above
result = has_close_elements([1.0, float('inf'), 3.0], 0.5)
assert result == False, "The function should return False when the list 
contains infinity, as infinity cannot be close to any finite number."
\end{verbatim}

In this test script, we are testing the function with lists that include \texttt{NaN} and \texttt{inf}. The function should ideally handle these cases gracefully, but due to the nature of these special floating-point values, the current implementation may not behave as expected. The assertions are designed to fail if the function does not correctly handle these cases.

\vspace{0.5cm} \hrule \vspace{0.5cm}

%==============================================================================
\noindent\textbf{\Large [L] Live Execution + Judge}
%==============================================================================

\textbf{Judge Prompt:}
\begin{verbatim}
You are an impartial Judge.

## Problem
[problem code]

## Proposed Test Case
[test script from Red Team]

## Task
Determine if this test case is VALID and CORRECT for the given problem.
1. Does the input satisfy all constraints (e.g. range, type, format)?
2. Is the asserted output (if any) logically correct?
3. Is it a fair test?

If the test case is valid, output: **VALID**
If it violates constraints or expects wrong output, output: **INVALID**
\end{verbatim}

\textbf{Judge Response:}

To determine if the proposed test cases are valid and correct for the given problem, let's analyze each aspect:

\textbf{1. Input Constraints:} The function \texttt{has\_close\_elements} accepts a list of floats and a float threshold. The proposed test cases use lists containing \texttt{float('nan')} and \texttt{float('inf')}, which are valid float values in Python. Therefore, the input satisfies the type constraints.

\textbf{2. Logical Correctness of the Asserted Output:}
\begin{itemize}
    \item For the test case with \texttt{float('nan')}: In Python, any comparison with \texttt{NaN} (Not a Number) is always false, including equality and inequality checks. Therefore, the presence of \texttt{NaN} should not affect the determination of whether any two numbers are closer than the threshold. The function should return \texttt{False} as no valid comparisons can be made with \texttt{NaN}.
    \item For the test case with \texttt{float('inf')}: Infinity is not a finite number, and it cannot be close to any finite number within a finite threshold. Therefore, the function should return \texttt{False} as no finite number can be within a finite threshold of infinity.
\end{itemize}

\textbf{3. Fairness of the Test:} The test cases are fair as they check the function's behavior with special float values (\texttt{NaN} and \texttt{Infinity}), which are edge cases that can occur in floating-point computations.

Given these points, the proposed test cases are logically correct and do not violate any constraints. They are valid tests for the function's behavior with special float values.

\textbf{Verdict: VALID}

\vspace{0.3cm}

\textbf{Live Execution Result:}
\begin{verbatim}
Executing test script with generated code...

def _concrete_verify_script(code, test_script, timeout=5):
    full_code = f"""
    {code}
    {test_script}
    """
    try:
        exec(full_code, exec_globals)
        return "PASS"
    except AssertionError:
        return "FAIL_ASSERT"
    except Exception:
        return "FAIL_CRASH"

Result: PASS (no assertion errors, no crashes)
\end{verbatim}

\vspace{0.5cm} \hrule \vspace{0.5cm}

%==============================================================================
\noindent\textbf{\Large Final Result}
%==============================================================================

\begin{center}
\fbox{\Large\textbf{PASSED after self-check}}
\end{center}

\vspace{3mm}

\begin{tabular}{ll}
\textbf{Total API Calls:} & 6 \\
\textbf{Total Time:} & 20.8 seconds \\
\textbf{Result:} & Solved (1/1 = 100\%) \\
\end{tabular}

\vspace{5mm}

The generated code correctly handles all edge cases including \texttt{NaN} and infinity without requiring any fixes. This demonstrates how SolidCoder's Live Execution component grounds verification in runtime reality rather than relying on mental simulation alone. The Judge correctly validated the Red Team's test cases as fair, and the Live Execution confirmed the code's robustness.

\end{tcolorbox}

\end{document}
